%% PhD Thesis Proposal
%% Knox, 2026.

\documentclass[english,11pt]{extarticle}
\usepackage[T1]{fontenc}
\usepackage[utf8]{inputenc}
\usepackage{geometry}
\geometry{verbose,tmargin=1in,bmargin=1in,lmargin=1in,rmargin=1in}
\usepackage{float}
\usepackage{booktabs}
\usepackage{tabularx}
\usepackage{amsmath}
\usepackage{amsthm}
\usepackage{amssymb}
\usepackage{bm}
\usepackage{graphicx}
\usepackage{units}
\usepackage{enumitem}
\usepackage{microtype}          % typographic polish: char protrusion, expansion
\usepackage{xcolor}
\usepackage{comment}
\usepackage{tcolorbox}          % highlighted result boxes
\tcbuselibrary{skins,breakable}
\usepackage{fancyhdr}           % running header with the short title
\usepackage[textsize=footnotesize]{todonotes}   % margin \todo{} notes; width follows \marginparwidth
\usepackage{authblk}

\makeatletter

\numberwithin{equation}{section}
\numberwithin{figure}{section}
\numberwithin{table}{section}

% ---- Theorem environments ----
\@ifundefined{thm}{%
  \theoremstyle{plain}%
  \newtheorem{thm}{Theorem}[section]%
}{}
\@ifundefined{prop}{\newtheorem{prop}[thm]{Proposition}}{}
\@ifundefined{lem}{\newtheorem{lem}[thm]{Lemma}}{}
\@ifundefined{cor}{\newtheorem{cor}[thm]{Corollary}}{}
\@ifundefined{conj}{\newtheorem{conj}[thm]{Conjecture}}{}
\@ifundefined{defn}{%
  \theoremstyle{definition}%
  \newtheorem{defn}[thm]{Definition}%
}{}
\@ifundefined{assum}{%
  \theoremstyle{definition}%
  \newtheorem{assum}[thm]{Assumption}%
}{}
\@ifundefined{example}{\newtheorem{example}[thm]{Example}}{}
\@ifundefined{remark}{%
  \theoremstyle{remark}%
  \newtheorem{remark}[thm]{Remark}%
}{}

% ---- Convenient macros ----
\newcommand{\R}{\mathbb{R}}
\newcommand{\C}{\mathbb{C}}
\newcommand{\N}{\mathbb{N}}
\newcommand{\E}{\mathbb{E}}
\newcommand{\Prob}{\mathbb{P}}
\newcommand{\ip}[2]{\langle #1, #2 \rangle}
\newcommand{\norm}[1]{\left\| #1 \right\|}
\newcommand{\abs}[1]{\left| #1 \right|}
\newcommand{\T}{^{\mathsf{T}}}          % transpose
\newcommand{\dd}{\,\mathrm{d}}          % upright differential
\newcommand{\D}[2]{\frac{\mathrm{d}#1}{\mathrm{d}#2}}     % total derivative
\newcommand{\pd}[2]{\frac{\partial #1}{\partial #2}}      % partial derivative

% ---- Program-wide notation ----
\newcommand{\Rset}{\mathcal{R}}          % reachable set
\newcommand{\Mplus}{\mathcal{M}_+}       % cone of nonnegative Radon measures
\newcommand{\Xset}{\mathcal{X}}          % state constraint set
\newcommand{\Uset}{\mathcal{U}}          % control constraint set
\newcommand{\Lie}{\mathcal{L}}           % Lie derivative / generator
\newcommand{\modes}{\mathcal{Q}}         % mode alphabet (Paper 1, Section 2.2)
\newcommand{\guard}{\mathcal{S}}         % guard surface
\newcommand{\reset}{\mathcal{J}}         % reset map

% The Henrion-Korda dual pair is conventionally (v, w). Both letters are
% already body-frame velocity components in the state vector, and the
% certificate is a function OF that state, so the two would collide inside a
% single displayed equation. Capital Greek instead: distinct from latitude
% \phi and heading \psi, which are lowercase. Named here so the choice is
% changed in one place rather than hunted through the papers.
\newcommand{\cert}{\Phi}               % Liouville supersolution, cert(t,x)
\newcommand{\certT}{\Psi}              % terminal majorant, certT(x)

\makeatother

% ---- Highlighted-result box ----
\newtcolorbox{keybox}[1][]{colback=blue!4,colframe=blue!45!black,
  boxrule=0.6pt,arc=2pt,left=6pt,right=6pt,top=4pt,bottom=4pt,#1}

% ---- Development annotations ----
%  Three coloured boxes carrying the drafting ledger inside the source:
%    decide -- a choice made, recorded so it is not silently re-opened
%    gap    -- an obligation not yet discharged
%    idea   -- a line of argument worth keeping, not yet load-bearing
%  They are deliberately conspicuous. Set \draftnotesfalse below (or pass
%  \PassOptionsToPackage before \input{preamble}) to compile a submission
%  copy in which all three vanish, exactly as [disable]{todonotes} removes
%  the \todo list --- the boxes must not survive into a submitted draft,
%  and a single switch is the only reliable way to be sure of that.
\newif\ifdraftnotes
\draftnotestrue

\newtcolorbox{decide@box}[1][]{colback=orange!4,colframe=orange!45!black,
  boxrule=0.6pt,arc=2pt,left=6pt,right=6pt,top=4pt,bottom=4pt,#1}
\newtcolorbox{gap@box}[1][]{colback=red!4,colframe=red!50!black,
  boxrule=0.6pt,arc=2pt,left=6pt,right=6pt,top=4pt,bottom=4pt,#1}
\newtcolorbox{idea@box}[1][]{colback=green!4,colframe=green!35!black,
  boxrule=0.6pt,arc=2pt,left=6pt,right=6pt,top=4pt,bottom=4pt,#1}

\ifdraftnotes
  \newenvironment{decide}{\begin{decide@box}}{\end{decide@box}}
  \newenvironment{gap}{\begin{gap@box}}{\end{gap@box}}
  \newenvironment{idea}{\begin{idea@box}}{\end{idea@box}}
\else
  \excludecomment{decide}
  \excludecomment{gap}
  \excludecomment{idea}
\fi

% ---- Theorem-environment aliases ----
%  The drafts write \begin{assumption} and \begin{lemma} in full; the
%  counters above are declared as assum and lem. Alias rather than rename,
%  so neither spelling is a compile error and no draft has to be rewritten
%  to match a decision about abbreviation.
\makeatletter
\@ifundefined{assumption}{%
  \newenvironment{assumption}[1][]{\begin{assum}[#1]}{\end{assum}}}{}
\@ifundefined{lemma}{%
  \newenvironment{lemma}[1][]{\begin{lem}[#1]}{\end{lem}}}{}
\makeatother

\usepackage{babel}

% ---- Cross-references and links (load late: hyperref then cleveref) ----
\usepackage[hidelinks,colorlinks=true,
  linkcolor=blue!55!black,citecolor=blue!55!black,urlcolor=blue!55!black]{hyperref}

% ---- Bibliography ----
%  biblatex, not the plain BibTeX \bibliography. Loaded after hyperref and
%  before cleveref, which is the order the biblatex manual asks for.
%
%  style=numeric keeps the bare-number in-text appearance the papers have
%  always used. sorting=none orders the reference list by first citation,
%  which reads better than alphabetical in a proposal.
%
%  sortcites is the one non-obvious option. It sorts the keys WITHIN each
%  citation so the numeric style can collapse consecutive labels into a
%  range -- without it a five-key \parencite renders as five separate
%  brackets. Note this is NOT natbib's sort&compress, which biblatex does
%  not provide: the biblatex natbib compatibility layer (blx-natbib.def)
%  defines \citep and \citet but nothing for compression, and an unknown
%  package key only warns, so a wrong setting here fails silently.
%
%  maxcitenames=2 keeps a multi-key citation from printing every author;
%  maxbibnames=99 with giveninits stops the reference list collapsing long
%  author lists to "et al." in a way that would not match the full lists
%  written out in shared.bib (see emmert2021nrlmsis, hersbach2020era5).
%
%  The per-entry reading notes in the bibliography live in the biblatex
%  `annotation' field. They are the most valuable content in that file and are
%  deliberately NOT printed. Note that `annotation' is a DATA field, not an
%  option of \printbibliography -- passing annotation=false to \printbibliography
%  is an xkeyval error, not a way of suppressing it. Blanking the field's format
%  is the mechanism, and it is the right one: the notes survive in the .bib and
%  in the .bbl for the author's use, and never reach the page. This must come
%  AFTER the \usepackage line below, since \DeclareFieldFormat is defined there.
\usepackage[backend=biber,style=numeric,sorting=none,sortcites,
  maxbibnames=99,giveninits=true,maxcitenames=2]{biblatex}

% biblatex requires csquotes whenever babel is loaded, and warns without it
% ("'babel/polyglossia' detected but 'csquotes' missing"). It supplies the
% quote punctuation the bibliography macros expect. Not currently used by any
% entry, but the dependency is unconditional.
\usepackage{csquotes}

\DeclareFieldFormat{annotation}{}

% The thesis bibliography. \addbibresource must be issued in the preamble, and
% this preamble is \input from a directory one level below manuscript/, so the
% path is written relative to that directory, not to this file.
\addbibresource{../thesis.bib}

\usepackage[nameinlink,capitalise]{cleveref}

% ---- Running header: each paper calls \runninghead{Short Title} ----
\newcommand{\runninghead}[1]{%
  \pagestyle{fancy}%
  \fancyhf{}%
  \lhead{\small #1}%
  \rhead{\small\thepage}%
  \renewcommand{\headrulewidth}{0.3pt}%
}

\runninghead{Certified analytical bounds for hypersonic reachable sets}
\hypersetup{
  pdftitle={Certified analytical bounds for hypersonic reachable sets},
  pdfauthor={Ethan Knox},
  pdfsubject={Preliminary PhD research proposal},
}

% ---- Summary-mode switch --------------------------------------------------
%  `make proposal`         -> full copy (with the three appendices).
%  `make proposal-summary` -> body-only copy for cold emails; it predefines
%  \summarymode, which drops the appendices and neutralises the parenthetical
%  pointers into them (\apxref / \apxeqref below), so the body stays clean
%  and no cross-reference dangles. Every pointer from the body into an
%  appendix must go through one of those two macros for that to hold.
\newif\ifsummary
\ifdefined\summarymode\summarytrue\fi
\ifsummary
  \newcommand{\apxref}[1]{}%
  \newcommand{\apxeqref}[1]{}%
  % Compress the reference list for the cold-email copy: 3 authors then "et al."
  \ExecuteBibliographyOptions{maxbibnames=3}%
\else
  \newcommand{\apxref}[1]{~(\cref{#1})}%
  \newcommand{\apxeqref}[1]{~(Eq.~\ref{#1})}%
\fi

\begin{document}

\title{\textbf{Certified analytical bounds for hypersonic reachable sets}\\[0.4em]
  {\large\normalfont\textit{A measure-theoretic profile decomposition in
  validated interval arithmetic}}}
\author{\textbf{Ethan Knox}}
\affil{\small\textit{Independent Researcher, Champaign, IL}}
\date{Preliminary PhD Research Proposal \\[0.5em] \today}

\maketitle

\begin{abstract}
The reachable set of a nonlinear control system can be embedded in a space
of occupation measures, where the dynamics become a linear constraint and
the reachable set is the union of the supports of the admissible terminal
measures. This research asks what that embedding yields under two
departures from the standard program: when the vector field is not a
function but a set-valued map, enclosed in validated interval arithmetic;
and when the measure family loses compactness along the fast-time axis in a
singular limit. On the state space nothing is lost---the family is
tight---but its weak limit forgets every event that lasts a vanishing time,
and a reachable set is a support, not an integral. The target theorem is a
profile decomposition in the manner of Lions' concentration-compactness
alternative, taken on the energy dissipation along the rescaled time axis:
a finite-dimensional skeleton, one boundary-layer bubble per event, and a
residual that is small in a weaker norm, with an explicit bound of
order~$\sqrt{\eps}$ as the aim. Almost-orthogonality of the scaling groups,
in the sense of the Bahouri--G\'erard construction and under bounded
adversarial uncertainty, is what would make the decomposition a bound
rather than an ansatz. The constants that size it are required in closed
form, derived from the physical constraint set and the aerodynamic interval
widths rather than from a numerical solver. Hypersonic entry supplies the
instance: a six-degree-of-freedom vehicle whose thin-atmosphere limit
furnishes the singular perturbation, whose equilibrium glide and
exo-atmospheric arcs furnish the skeleton, and whose atmospheric passes
furnish the bubbles. The aim is a certificate rather than a computation:
explicit algebraic bounds, valid uniformly over the entire admissible
family, in the minimal physical state space.
\end{abstract}
\pagebreak

% \listoftodos[Open items]
% \pagebreak

\ifsummary\else
\tableofcontents
\pagebreak
\fi

%==========================================================================
\section{The analytical problem}
%==========================================================================
The object to be bounded is the reachable set of an atmospheric entry
vehicle---every state it can be driven to under admissible control---which
is what safety and threat analyses need enclosed. The flight-mechanics
specifics are deferred: \cref{sec:application} states why entry is the
right instance and what its state space looks like%
\ifsummary
, and the coupled dynamics---frames, state vector, translational and
rotational equations, the aerodynamic bridge from attitude to force, and
the glide condition---are carried, together with sketches of the limit
argument and of the constants it rests on, in companion appendices omitted
from this copy%
\else
, \cref{apx:dyn} carries the coupled dynamics in full, and
\cref{apx:time,apx:residual} sketch the limit argument and the constants it
rests on%
\fi
. What matters first is that the
problem is an analytical one before it is an aerospace one, and that it
fails for two independent reasons---the dimension of the state, and the
provenance of the guarantee.

\subsection{Reachable sets as occupation measures}\label{sec:embedding}
Embed the hybrid flight dynamics into a Banach space of occupation
measures. The nonlinear flight mechanics $\dot x = f_q(x,c,t)$ become
\emph{linear}: the Liouville equation is a linear constraint on a measure
$\mu$ over $[0,T]\times\calX\times\calU$ and its terminal marginal
$\mu_T$, and the reachable set at time~$T$ is the union of
$\operatorname{supp}\mu_T$ over the measures that satisfy it. No single
measure bounds it. The outer bound comes from the dual: a pair of
functions $(\cert,\certT)$, the first non-decreasing along every admissible
trajectory and at least one on the initial set, the second majorizing it at
time~$T$, encloses the reachable set in the superlevel set
$\{\certT\ge 1\}$~\cite{henrion2014convex,korda2014convex}. So far
this is the occupation-measure program of Lasserre, Henrion and
collaborators~\cite{lasserre2008nonlinear,vinter1993convex}. The embedding buys linearity
at the cost of uniqueness, and the way it is recovered is standard: a
trajectory family satisfying the Liouville constraint lies in the convex
hull of the generalized curves, so measures correspond to convex mixtures
of trajectories rather than to one of them~\cite{henrion2014convex,ball1989version}.

The same relaxation admits the interval envelope. Because the true
aerodynamics are transcendental and known only to within an envelope, they
are enclosed in intervals, and a vector field selected from an interval is
a set-valued map; the resulting differential inclusion is the object the
measure must describe, not the nominal
field~\cite{filippov1988differential,aubin1984differential}. Written with
the selection as an additional box-valued input, the inclusion is a control
system with a bounded disturbance, and at fixed $\eps$ its Liouville
constraint is the standard one. What is not standard is carrying the outer
approximation through the singular limit of \cref{sec:profile}, uniformly
over the selections. Occupation-measure theory for inclusions of exactly
this kind already exists, and it is the nearest prior art: it defines an
occupational invariant measure without assuming solution uniqueness, which
is the same move made here~\cite{artstein1999invariant}. What it assumes is
worth stating plainly, because it locates where the entry problem leaves
it. The linear growth condition on the field serves only to exclude
finite-time blow-up, and is automatic on a bounded corridor. The
singular-perturbation theorem, however, assumes a uniform bound on the
fast state along the whole
$\eps$-family~\cite{artstein1999invariant,artstein2000value}, and the
natural fast state of entry is the altitude measured in scale heights: on a
corridor of fixed altitude extent it ranges over an interval of length
proportional to $1/\eps$, the aerodynamic terms grow exponentially in it,
and a skip carries it out of every bounded set. The hypothesis fails
exactly at the events of interest.

\begin{openitem}{Interval embedding, uniform through the singular limit}
\emph{Question.} Does the occupation-measure outer approximation of the
reachable set survive, with constants independent of $\eps$, when the
vector field is an interval inclusion $\dot x \in \calF_\eps(x,c,t)$?

\emph{To establish.} At fixed $\eps$ the inclusion is a control system with
a box-valued disturbance and its Liouville constraint is standard. The work
is the uniformity: a certified upper and lower envelope of the atmospheric
density, which fixes $\calF_\eps$ concretely\apxref{sec:density}, and a
proof that every trajectory of every selection is represented by a feasible
measure whose bounds on mass, support and dissipation do not degrade as
$\eps\to 0$.

\emph{Success.} The embedding stated for the full inclusion, with each
hypothesis either verified on the entry corridor or named as an assumption.
\end{openitem}

\subsection{The provenance of the guarantee}\label{sec:provenance}
A polynomial barrier whose coefficients come from an interior-point solver
has analytical \emph{form} but inherits the solver's
accuracy~\cite{prajna2007framework}, and the documented failure modes of
numerical semidefinite solvers include the loss of strict feasibility on
which the repair schemes depend~\cite{roux2018validating}. A solver's
output can be made into a certificate: rounded to rational coefficients and
re-verified exactly, it is a proof~\cite{peyrl2008computing}. The obstacle
is scale, not principle. The rounding needs a strictly feasible Gram
matrix, which a tight bound does not have, and exact rational certificates
carry a documented bit-complexity whose exponent grows with the number of
variables~\cite{kaltofen2012exact,magron2021exact}---a cost the
thirteen-state field of \cref{sec:application} does not invite as a main
line.

The program therefore fixes an order of evidence. The certificate itself is
a closed-form bound: the constants of \cref{sec:profile}, derived from the
decomposition and evaluated over the corridor in validated interval
arithmetic. Behind it stand symbolic checks of the derivations those
constants rest on, run as part of the build against the same definition of
the dynamics that is evaluated numerically. A moment-sum-of-squares
computation on a reduced model is the third tier: an independent
cross-check, to which the exact-certification machinery above belongs, and
never the source of the guarantee.

Interval arithmetic is the setting in which a guarantee has meant, for
forty years, an enclosure rather than an approximation. It is rigorous in
floating point: every operation is rounded outward, so the computed
interval contains the exact result, and nothing is gained for soundness by
exact rational arithmetic, which cannot in any case evaluate an exponential
or a sine~\cite[Ch.~3]{moore2009introduction}. Two distinct losses have to
be managed, and they are often run together. The first is the
\emph{dependency problem}. Interval arithmetic is inclusion
isotonic---the property that makes an enclosure valid at all---but not
every law of real arithmetic survives the move to intervals: distributivity
weakens to subdistributivity, and repeated occurrences of a variable are
treated as independent~\cite[Chs.~4--5]{moore2009introduction}. The second
is the \emph{wrapping effect}: the image of a box under a flow is not a
box, and re-enclosing it at every step admits points that no trajectory
reaches~\cite[Ch.~10]{moore2009introduction}. Refining the step does not
help. Under interval Euler steps of size~$h$, re-boxed at each, a rigid
rotation inflates a box by $(1+h)^{2\pi/h}$ per revolution, which
\emph{increases} to $e^{2\pi}\approx 535$ as $h\to 0$; what controls
wrapping is the
representation of the set---Lohner's QR factorization, zonotopes, Taylor
models~\cite{nedialkov1999validated,makino2003taylor}. Applied interval
analysis supplies the machinery the pipeline is built from, and it is
deliberately conservative: certified set inversion, contractors, and
validated integration~\cite{jaulin2001applied}.

\begin{openitem}{Validated computational pipeline}
\emph{Question.} Can the closed-form constants, and the pass enclosures
they do not cover, be evaluated over the whole corridor with every rounding
error accounted for?

\emph{To build.} An interval evaluator for the transcendental field and its
Jacobian with outward rounding, contractors for the corridor constraints,
and a validated integrator for the passes, extending the ball-arithmetic
prototype described in \cref{sec:status}.

\emph{Success.} Every inequality the certificate asserts is discharged by a
program that returns either a proof or an explicit failure to decide at the
working precision, and the dynamics it evaluates are the dynamics that were
checked symbolically.
\end{openitem}


%==========================================================================
\section{The singular limit and its profile decomposition}
\label{sec:profile}
%==========================================================================
This section carries the new content. The instance supplies a small
parameter---in the entry problem the thin-atmosphere ratio
$\eps_{\mathrm{atm}} = H_s/R_e \to 0$, where $H_s$ is the atmospheric
scale height and $R_e$ the planetary radius, written $\eps$ where the other
small parameters are not in play---and the question is what the
occupation-measure family does in that limit.

\subsection{Loss of compactness along the time axis and the defect measure}
On the corridor nothing is lost in the usual sense. The set
$[0,T]\times\calX\times\calU$ is compact, so the occupation measures---which
live in the cone $\Mplus([0,T]\times\calX\times\calU)$ of nonnegative
Radon measures, topologized weakly-$*$ as the dual of
$C([0,T]\times\calX\times\calU)$---form a tight family and converge along
subsequences. The limit is nevertheless the wrong object, because it
forgets the atmospheric passes. A pass lasts a time of order
$\sqrt{\eps_{\mathrm{atm}}}$, the phugoid-to-secular time ratio (which, as
in the classical Chapman--Loh variables, enters as
$\sqrt{\eps_{\mathrm{atm}}}$, not
$\eps_{\mathrm{atm}}$~\cite{chapman1959approximate,loh1968reentry}). It
therefore carries occupation mass of that order, and none in the limit,
while its share of what distinguishes one trajectory from another---the
energy dissipated, the heat absorbed, the turn of the velocity
vector---need not vanish with it. A reachable set is a support, not an
integral, and what has no mass still counts. That is the gap between this
problem and the limit theory of value
functions~\cite{artstein2000value}: an integral survives the loss of a null
set, and a support does not. The loss has to be read where it occurs, along
the fast-time axis $\tau = t/\sqrt{\eps_{\mathrm{atm}}}$, on which the
horizon $[0,T/\sqrt{\eps_{\mathrm{atm}}}]$ exhausts the half-line and
translations act non-compactly.

Between passes the vehicle is outside the sensible atmosphere, on an arc
governed by gravity alone; once captured, it stays near the
equilibrium-glide balance $\Lambda = 1$. In the 3-DOF (point-mass)
approximation this is the scalar condition
$L\cos\sigma/[m(g - V^2/r)] = 1$; in the full 6-DOF formulation the lift
component $L\cos\sigma$ is not an independent input but is determined by
the attitude quaternion through the body-to-LVLH
rotation\apxref{sec:aero-force,sec:glide-manifold}, so the balance is
a single scalar constraint on the 13-dimensional translational-rotational
state, and the slow set is its intersection with the trim and
quasi-static-deflection conditions of the two faster groups---one
constraint set for each. The balance is taken with planetary rotation and
oblate gravity retained\apxeqref{eq:glide-full}, and the two do not enter
on equal terms: flying east at entry speed the Coriolis term is tens of
percent of the net acceleration $g - V^2/r$ that the lift has to balance,
while the $J_2$ terms are a percent of it at most\apxref{rem:j2-glide}. The
scalar $\Lambda = 1$ is the reduction under spherical gravity, no rotation
and small flight-path angle.

The loss is read on a \emph{defect measure}, and the object has to be chosen
so that the passes carry mass. The part of the occupation-measure limit
that lies off the glide cannot serve, since the passes have none. The
energy dissipation can. Along every trajectory of the inclusion the
specific mechanical energy in the rotating frame,
$\mathcal{E} = \tfrac12 V^2 + \Phi_g - \tfrac12\omega_E^2 r^2\cos^2\phi$,
changes at the rate $\dot{\mathcal{E}} = -P$, where $P = DV/m \ge 0$ is the
drag power: neither the lift nor the Coriolis force does
work\apxref{lem:energy}. The total dissipation
$\int_0^T P\dd t \le \mathcal{E}_{\max} - \mathcal{E}_{\min}$ is therefore
bounded uniformly over controls, over selections of the interval field,
and over $\eps$. Let
\begin{equation}\label{eq:defect}
  \nu_\eps(\tau)
  \;=\; \frac{\sqrt{\eps}\;P\bigl(\sqrt{\eps}\,\tau\bigr)}
             {\int_0^T P\dd t}
\end{equation}
be the dissipation density on the fast-time axis, normalized to unit mass
and extended by zero. Along any sequence $\eps_n\to 0$, with a trajectory
chosen for each, these are nonnegative densities of fixed mass on the
line, which is exactly the hypothesis of Lions' concentration-compactness
lemma~\cite[Lemma~I.1]{lions1984concentration1}, and here all three of its
alternatives mean something:
\begin{itemize}[itemsep=1pt]
  \item \emph{vanishing}---no window of fixed fast-time width holds a fixed
    share of the dissipation: the glide, on which the drag is bounded and
    such a window holds a share of order $\sqrt{\eps}$;
  \item \emph{compactness}, up to translation---one pass carries the
    dissipation;
  \item \emph{dichotomy}---the dissipation splits between pieces whose
    fast-time separation diverges: two passes a positive time apart, or a
    pass followed by a glide.
\end{itemize}
Translations along the fast-time axis therefore separate the passes, and it
is the dilation half of the principle---the limit-case
modification~\cite{lions1985limit1}---that resolves each one at its own
width, since a pass shorter than $\sqrt{\eps}$ is a Dirac mass at that
scale. The trichotomy is Lions' lemma applied as stated. That the vanishing
part is the glide and the concentrated pieces are the passes is asymptotics
still to be argued, from the inner and outer expansions of the classical
theory~\cite{shi1971skip}, and one step of the argument is
open\apxref{apx:time}.

\subsection{Skeleton, bubbles, residual}
Rather than restore compactness, we extract the profile decomposition that
iterating the dichotomy hands us:
\begin{itemize}
  \item a finite-dimensional \emph{skeleton} $S$---what survives the limit:
    the exo-atmospheric arcs between passes, which are conics up to the
    oblateness correction, followed by the equilibrium glide, a graph over
    $(r,V)$ in the corridor subspace once the balance is solved for the
    vertical lift fraction. For the glide a rigorous characterization with
    quantitative accuracy estimates already
    exists~\cite{lu2006asymptotic}, and it treats the glide as a
    \emph{regular} perturbation: the exact solution is not attracted to the
    zeroth-order one, a point that returns in \cref{sec:bound};

  \item one \emph{boundary-layer bubble} per atmospheric pass, the analogue
    of a concentrating profile: the pass rescaled to its own width, whose
    net effect on the slow coordinates is a compact impulse set $\calI_j$,
    with orthogonality read off the scale-and-core separation condition
    that governs concentrating profiles
    generally~\cite{gerard1998description}. The decomposition of skip entry
    into an aerodynamic inner region and a gravitational outer one is
    classical~\cite{shi1971skip,chapman1959approximate}. What is not
    classical is a statement uniform over a family of trajectories, with an
    explicit remainder at fixed $\eps$;

  \item a \emph{residual}: what the finitely many profiles do not capture.
    In a profile decomposition the remainder is not compact and is not
    small in the conserved norm; it is small in a weaker one, and only as
    the number of extracted profiles
    grows~\cite{bahouri1999high,solimini1995note}. Here the weaker norm is
    the largest share of dissipation held by any fast-time window: on a
    glide the residual carries most of the dissipation, and almost none of
    it in any one window. The quantitative bound the program needs is a
    different statement, a distance to the skeleton in state space, and it
    is taken up in \cref{sec:bound}.
\end{itemize}

The number of bubbles is bounded \emph{a~priori}, and the bound comes from
the same monotone quantity. A pass that dissipates at least $\delta$ of
specific energy can occur at most
\begin{equation}\label{eq:kmax}
  K_\delta \;\le\; \frac{\mathcal{E}_{\max}-\mathcal{E}_{\min}}{\delta}
\end{equation}
times, for every control and every selection of the interval field. Passes
below the threshold are assigned to the residual, which is the device of
the profile-decomposition literature: only finitely many profiles exceed
any given size. The integrated heat-load budget gives the count its
physical reading---a trajectory can only afford so many
skips~\cite{shi1971skip,eggers1958comparative}---but it routes through a
stagnation-point heating correlation, and the energy budget does
not\apxref{rem:sutton-graves-circularity}.

Enclosing a pass is where the wrapping effect has to be met. For linear
systems it can be avoided altogether: support-function methods propagate
each direction independently, so successive over-approximations do not
compound, and what remains is a discretization error of first order in the
time step~\cite{leguernic2010reachability}. Nothing comparable is available
for a nonlinear pass.

\begin{openitem}{Enclosures of the pass maps}
\emph{Question.} Can the net effect of one atmospheric pass on the slow
coordinates be enclosed tightly enough that a sequence of $K_\delta$
passes does not exhaust the corridor?

\emph{To prove and build.} Wrapping in a pass has a specific source.
Trajectories of a set enter the layer at different times, the set is
sheared along the flow, and re-enclosing the sheared set in a box at every
step admits points that no trajectory reaches. The remedy that fits the
decomposition is to enclose the \emph{pass map}---entry state to exit
state, in the rescaled time of the profile---as a single object, a
zonotope or Taylor-model enclosure of a map rather than a tube of boxes,
and to compose pass maps with the outer flow.

\emph{Success.} A bound on the growth of the enclosure per pass, explicit
in the interval widths, and its verification on a multi-pass corridor
against sampled trajectories.
\end{openitem}

\subsection{Almost-orthogonality of the scaling groups}
Three scaling groups act on the time axis. The atmospheric one dilates by
$\sqrt{\eps_{\mathrm{atm}}}$; those of the attitude and the elastic motion
dilate by $\eps_{\mathrm{att}} = T_{\mathrm{att}}/T_{\mathrm{slow}}$ and
$\eps_{\mathrm{ela}} = T_{\mathrm{ela}}/T_{\mathrm{slow}}$, the
short-period attitude and first structural periods measured against the
slow time $T_{\mathrm{slow}} = \sqrt{R_e/g}$. Their separation is the
Bahouri--G\'erard almost-orthogonality condition, read on the dissipation:
two profiles are orthogonal when their scales diverge or, at equal scale,
when their centers separate by a diverging number of
widths~\cite[(1.18)]{bahouri1999high}. Two things follow from it, and they
differ in kind. The dissipation is a measure, so it splits
\emph{additively} across orthogonal profiles; that is the analogue of the
energy identity of the wave-equation theory, and it costs nothing. The
state does not split additively, because the profiles interact through the
nonlinear dynamics: the state leaving one pass is the datum of the next
arc. What has to be proved is the analogue of the nonlinear statement of
that theory, that the contributions of separate profiles compose without
interacting, up to an error that is explicit and uniform over the family.
For passes separated in time this is a matching estimate between inner and
outer solutions. For the three groups it is scale separation where the
timescale bands are ordered; where two bands cross inside the corridor the
groups are not orthogonal there, and the decomposition has to carry one
joint profile on the coupled subsystem rather than two. Establishing this
with the scales themselves known only to within the aerodynamic intervals
is the one genuinely new obligation, and it sits on the critical path.

\begin{openitem}{Almost-orthogonality under interval disturbances}
\emph{Question.} Do the scaling groups remain orthogonal, in the sense of
the Bahouri--G\'erard condition, for every selection of the aerodynamic
interval field---and where they do not, what replaces the decomposition?

\emph{To prove.} A non-interaction theorem: for profiles separated in scale
or in time, the map from entry data to exit data factors through the
individual pass maps and the outer flow, with an explicit remainder,
uniformly over selections. The interval widths enter twice: as a
disturbance within each profile, and through the periods $T_{\mathrm{att}}$
and $T_{\mathrm{ela}}$, which depend on the uncertain coefficients, so that
the ordering of the bands is itself a hypothesis to be verified over the
corridor. Where two bands cross, the statement is for a joint profile.

\emph{Success.} The theorem for the atmospheric and attitude groups with
the remainder in closed form; the elastic group either included or
enclosed as a widened interval, as in the fallback of \cref{sec:plan}.
\end{openitem}

The argument worth transporting is the constructive one: splitting a
bounded sequence into an oscillatory and a singular part at each scale
forces those scales to be orthogonal, and iterating extracts a pairwise
orthogonal family, so orthogonality is produced by the construction rather
than assumed by it~\cite{bahouri1999high}.

\subsection{The certified bound}\label{sec:bound}
The target statement is
\begin{equation}\label{eq:target}
  \calR_\eps(T)
  \;\subseteq\;
  \biggl(
    \underbrace{\Phi_S(\Theta)}_{\text{skeleton}}
    \;\oplus\;
    \underbrace{\bigoplus_{j=1}^{K_\delta}\calI_j}_{\text{bubbles}}
  \biggr)
  \;\oplus\;
  \underbrace{\B_\varrho\bigl(C_{\calR}\sqrt{\eps}\bigr)}_{\text{residual}},
\end{equation}
uniformly over the admissible family generated by the uncertainty set, not
along a nominal trajectory. Here $\calR_\eps(T)$ is the set of slow states
reachable at time~$T$ from the initial set, under every admissible control
and every selection of the interval field. $\Theta$ is the compact set of
skeleton parameters---initial data and slow controls---and
$\Phi_S\colon\Theta\to\calX$ the endpoint map of the skeleton flow.
$\calI_j$ is the set of displacements of the time-$T$ endpoint
attributable to pass~$j$: its impulse set, carried to time~$T$ by the
skeleton flow, over all states from which the pass can be entered. The sum
$\oplus$ is the Minkowski sum, and $\B_\varrho(a)$ is the ball of
radius~$a$ in the corridor metric~$\varrho$, a weighted supremum norm in
which altitude is measured in scale heights and each remaining slow
coordinate in units of its corridor width. Passes compose in sequence, not
in parallel, and the sum is the telescoped form of that composition: the
endpoint of a trajectory with $K$ passes differs from the skeleton endpoint
by a sum of $K$ displacements, one from each $\calI_j$. The sum forgets
that successive displacements are correlated, which makes the bound
conservative and not unsound.

The residual constant $C_{\calR}$ and the bubble count $K_\delta$ are to be
derived as explicit, closed-form expressions in the corridor bounds and
the aerodynamic interval widths, independent of any numerical
optimization. Both exist in first form. The count is \cref{eq:kmax}. On
the glide, with $L/D$ the lift-to-drag ratio and $\sigma$ the bank angle,
\begin{equation}\label{eq:CR-body}
  C_{\calR}
  \;=\; \frac{2\sqrt{g R_e}}
             {V\,(L/D)\cos\sigma\,\sqrt{1 - V^2/(g r)}} ,
\end{equation}
which bounds the computed distance to the glide to within a few percent
over the range of lift-to-drag ratio, speed and bank
tested\apxref{apx:residual}. The same analysis shows what the formula does
not yet cover, and the limitation is structural. $C_{\calR}\sqrt{\eps}$ is
twice the damping ratio of the phugoid, the oscillation of altitude against
flight-path angle about the glide. Its damping \emph{rate},
$g/[V(L/D)\cos\sigma]$, does not depend on $\eps$, while its frequency
grows like $\eps^{-1/2}$. In the limit the motion normal to the glide is an
undamped oscillation on its own timescale: the glide balance is not a
normally hyperbolic slow manifold in the sense the geometric theory
requires~\cite{fenichel1979geometric}, and no boundary-layer estimate of
Tikhonov type brings a trajectory onto it. The constant above is the
ringing amplitude of a trajectory started level on the glide; a trajectory
leaving a pass rings with whatever amplitude the pass left it, for a time
comparable to the entry. The residual bound therefore has to carry that
amplitude, either as an explicit decaying term or by promoting the phugoid
amplitude to a slow coordinate of the skeleton and averaging over its
phase; and $C_{\calR}$ diverges as $V$ approaches circular speed, which is
where the glide gives way to arcs and passes.

In the spherical-LVLH coordinates the altitude and speed bounds of the
corridor are box constraints on $(r,V)$. The heating, dynamic-pressure and
load limits are not: each is a curve $h \ge h_{\min}(V)$. Each is, however,
monotone in altitude and in speed at fixed attitude, so contracting a box
against it returns the exact interval hull of the feasible part, and the
interval widths that enter $C_{\calR}$ and $K_\delta$ inherit that
tightness\apxref{sec:coord-rationale}. Requiring the constants in closed
form, rather than as the output of a solver, is what separates a
certificate from a computation in the sense of \cref{sec:provenance}.

\begin{openitem}{Closed-form constants}
\emph{Question.} What are $C_{\calR}$ and $K_\delta$ as functions of the
corridor bounds, the initial set and the aerodynamic interval widths?

\emph{To derive.} The count from the energy budget, with the threshold
$\delta$ chosen against the residual it leaves behind. The residual in
three parts: the lag behind the glide, which \cref{eq:CR-body} already
gives; the ringing left by a pass, which decays at a rate independent of
$\eps$ and has to be carried explicitly or absorbed into the skeleton by
averaging; and the matching error at each pass. Each constant is to hold
for every selection of the interval field, and its range of
validity---bounded away from circular speed---is part of the statement.

\emph{Success.} Formulas whose evaluation over the corridor in validated
interval arithmetic bounds the distance from sampled trajectories of the
full model to skeleton plus bubbles, with no fitted number.
\end{openitem}


%==========================================================================
\section{The application: hypersonic entry}\label{sec:application}
%==========================================================================
Entry is not an arbitrary illustration of the machinery above. It supplies
the singular perturbation, the skeleton and the concentrating events in one
system, and it is a problem the existing methods are known not to reach.

\subsection{Why existing methods do not reach it}
The direct method solves a Hamilton--Jacobi--Isaacs equation on a grid,
which is exponential in state dimension and capped in practice near four or
five states~\parencite{bansal2017hamilton,mitchell2005time}.
A rigid-body entry vehicle carries thirteen states,
\begin{equation}\label{eq:state-rigid}
  x = (r,\lambda,\phi,\, V,\gamma,\psi,\,
       q_0,q_1,q_2,q_3,\, p,q_b,r_b)
  \;\in\;\R^{13},
\end{equation}
and $17+2N_m$ once integrated heat load, wall temperature, ablation
recession, mass loss and $N_m$ aerothermoelastic modes are
coupled~\parencite{mcnamara2011aeroelastic}. The two counts use
``degree of freedom'' in its two customary senses, and it is worth
separating them: the vehicle is six-degree-of-freedom in the
configuration sense (three translational, three rotational), while
\cref{eq:state-rigid} has thirteen coordinates spanning a
twelve-dimensional state manifold once the unit-quaternion constraint
$\norm{\vvec{q}} = 1$ is imposed.

So the grid method stops before the real problem starts.

Set-propagation methods go further. Taylor-model and zonotope flowpipes
return rigorous outer bounds at ten to twenty states, and they are the
baseline this program has to be measured
against~\cite{chen2013flowstar,althoff2021set,makino2003taylor}. What they
return is an enclosure of one instance: a numerical set for one vehicle,
one initial box and one horizon, whose tightness tends to degrade with the
width of the initial set and with the length of weakly damped motion it has
to follow---and the phugoid is weakly damped for the whole of the glide.
They are the right tool inside a pass, where the horizon is short, and that
is how objective~3 uses them. They do not produce a bound whose dependence
on the corridor and on the interval widths is explicit, which is what a
statement uniform over a family requires.

What makes the measure-theoretic route viable is not that it evades the
dimension but that the bound it produces does not scale with it: the
skeleton is low-dimensional by construction, the bubbles are compact, and
the residual is a single ball.

\subsection{The state space and the corridor}
The translational state is carried in geocentric spherical position
$(r,\lambda,\phi)$ and wind-frame velocity $(V,\gamma,\psi)$, so that the
altitude and speed bounds of the entry corridor are box constraints on
states and the heating, dynamic-pressure and load limits are monotone in
the two coordinates they bind---the natural setting for interval
contraction. Attitude is a unit quaternion, free of Euler-angle
singularities. The chart has a price, which is stated rather than hidden:
the longitude and heading rates are singular at the geographic poles, so
the corridor is taken with $\abs{\phi}\le\phi_{\max}<\pi/2$, and a
trajectory class that overflies a pole needs a second chart or a Cartesian
formulation\apxref{sec:coord-rationale}.

Because interval arithmetic evaluates transcendentals natively, the
algebraic adjunctions that recast the system in polynomial
form~\cite{savageau1987recasting} are not needed on the main line; they
survive only inside the moment-sum-of-squares cross-check. The state space
therefore carries only the physical dimensions---spherical position,
wind-frame velocity, attitude quaternion, body rates, thermal, ablation,
and modal states---with no artificial embeddings, and the two auxiliary
states of the polynomial formulation, the density root
$y = \sqrt{\rho/\rho_0}$ and the wall-temperature surrogate, are
gone\apxref{sec:thermal}. This is a design decision rather than an open
problem, and its trade is stated once: transcendental dynamics in
coordinates native to the corridor and the glide, against polynomial
dynamics in which the corridor is a nonlinear function of the
state\apxref{sec:summary-table}.

\begin{openitem}{Independent verification}
\emph{Question.} Is the certificate right, by evidence that does not share
its derivation?

\emph{To build.} Three checks, in the order of \cref{sec:provenance}.
Symbolic verification of every derivation behind the constants, run with
the build. A moment-sum-of-squares outer approximation on a reduced
polynomial model, as a cross-check that shares neither the decomposition
nor the interval arithmetic. And a comparison on a common benchmark---a
capsule skip entry reconstructed from published data---against Taylor-model
flowpipes and against sampled trajectories of the full model, reported for
each coordinate group rather than as a single number.

\emph{Success.} Containment of every sample by the certified set, agreement
with the cross-check where both apply, and a stated tightness against the
baseline, coordinate by coordinate.
\end{openitem}


%==========================================================================
\section{Plan of work}\label{sec:plan}
%==========================================================================
Six concrete obligations carry the program; they are stated where each
arises above and gathered here, in order, as the plan of work.
\begin{enumerate}[leftmargin=*,itemsep=2pt]
  \item \emph{Interval embedding, uniform through the singular limit}: the
    outer approximation for the set-valued (interval) dynamics, with
    constants independent of $\eps$.
  \item \emph{Validated computational pipeline}: interval evaluation,
    contraction and integration that discharge the certificate with every
    rounding error accounted for.
  \item \emph{Enclosures of the pass maps}: tight enclosures carried through
    a sequence of atmospheric passes without premature blow-up.
  \item \emph{Almost-orthogonality under interval disturbances}:
    non-interaction of the scaling groups under bounded
    uncertainty---\textbf{the critical path}, on which the profile
    decomposition and the closed-form bound rest.
  \item \emph{Closed-form constants}: $C_{\calR}$ and $K_\delta$ from the
    physical corridor and the interval widths, free of any numerical
    solver.
  \item \emph{Independent verification}: symbolic checks of the
    derivations, a moment-sum-of-squares cross-check, and a benchmark
    against set-propagation tools and sampled trajectories.
\end{enumerate}
Objectives 1--2 establish the interval occupation-measure embedding, 3--4 the
profile decomposition, and 5--6 reduce it to closed form and check it;
objective 4 gates the rest.

\subsection{Where the work stands}\label{sec:status}
The computational layers of the program have working prototypes; the
central theorem is open. An implementation of the coupled model is public,
with a suite of
verification tasks whose failure criteria are stated before they
run.\footnote{\url{https://github.com/ethank5149/aether}} Three of its
results bear directly on the claims above. The residual scaling has been
measured rather than assumed: with the scale height varied over more than a
decade, the distance to the glide scales as $\eps^{0.49}$, and the closed
form \cref{eq:CR-body} bounds it from above to within about four percent.
The translational equations of motion, the energy balance behind
\cref{eq:kmax} and the linearization behind $C_{\calR}$ are reduced to
identities by computer algebra, against the same definition of the dynamics
that the code evaluates in floating point and in interval arithmetic, and
compiled numerical routines are generated from those expressions. And a validated-arithmetic prototype exists: enclosures
of the entry field and reachable tubes in ball arithmetic, with a certified
bound on the downrange still available to a capsule on a published skip
entry, holding for every bank history that keeps the trajectory inside a
stated corridor. None of this is the theorem. It is evidence that the
constants are computable and that the pipeline can be built, and the same
suite records the obstacle of \cref{sec:bound}, the weakly damped phugoid,
as a computed fact rather than a conjecture.

\subsection{Sequencing}
The dependency structure suggests three phases rather than six parallel
tracks. The first establishes the embedding---objective~1, together with the
certified envelope of the atmospheric density that it needs in order to be
stated on a concrete inclusion\apxref{sec:density}. The second attacks the
critical path: objective~4, which is where the program either works or does
not, and which should therefore be reached early enough that a negative
result is still recoverable. Objective~3 runs alongside it, since the
almost-orthogonality argument and the enclosure geometry constrain each
other. The third phase reduces the decomposition to closed
form---objective~5---completes the validated pipeline of objective~2, and
runs the symbolic and moment-sum-of-squares checks of objective~6 against
the result. The pipeline is the only deliverable that cannot be shortened
by a weaker theorem.

Two retreat positions are worth identifying now. If objective~4 resists a
proof for all three scaling groups simultaneously, the fallback is to
establish it for the atmospheric and attitude groups and to treat the
elastic group as a widened interval rather than a separate profile. That
weakens the residual constant without invalidating the decomposition. And
if the ringing left by a pass cannot be bounded in closed form, the phugoid
amplitude is carried as one more slow coordinate of the skeleton, which
costs a dimension and keeps the bound.

\ifsummary\else
\pagebreak

\appendix
\crefalias{section}{appendix}
\crefalias{subsection}{subappendix}
\crefalias{subsubsection}{subsubappendix}

% =====================================================================
%  Thesis proposal, Appendix: Coupled Entry Dynamics and Embedding
%  Owns: apx:dyn
% =====================================================================

\section{Coupled entry dynamics and embedding}
\label{apx:dyn}

%--------------------------------------------------------------------------
\subsection{Coordinate rationale}\label{sec:coord-rationale}
%--------------------------------------------------------------------------

The coordinate choice is dictated by two competing requirements of the
interval-arithmetic pipeline:
\begin{enumerate}
  \item \textbf{Corridor constraints must be box constraints.} The entry
    corridor---bounded altitude, bounded speed, heating and load-factor
    limits---defines the region whose reachable subset we certify. In
    coordinates where these constraints are box faces, the interval hull is
    the natural bounding object and no wrapping is introduced by
    intersecting with the corridor. Spherical position~$(r)$ and
    wind-frame speed~$(V)$ make the primary corridor constraints
    $h_{\min}\le r - R_e \le h_{\max}$ and $V_{\min}\le V\le V_{\max}$
    into box constraints on states.

  \item \textbf{The slow manifold must be native.} The equilibrium-glide
    condition $\Lambda = 1$ is a codimension-1 surface in the
    $(r,V,\gamma)$ subspace---equivalently, once solved for the flight-path
    angle, a graph over~$(r,V)$. The profile decomposition---skeleton, bubbles,
    residual---is organized around this surface. If the slow manifold is a
    derived nonlinear function of the states (as it would be in Cartesian
    position/velocity), the separation into skeleton and bubbles inherits
    unnecessary wrapping at every step.
\end{enumerate}

The translational state is therefore carried in spherical-LVLH coordinates
$(r,\lambda,\phi,V,\gamma,\psi)$. Attitude is carried as a unit quaternion
$\vvec{q} = (q_0,q_1,q_2,q_3)^\tr$ relating the body frame~$\mathcal{B}$ and
the ECEF frame~$\mathcal{E}$ (its direction-cosine matrix
$\mat{R}_{\mathcal{E}\mathcal{B}} = \mat{R}(\vvec{q})$ is given in
\cref{eq:DCM})---this avoids all Euler-angle singularities and keeps the
quaternion kinematics in their simplest form.

The longitude equation
$\dot\lambda = V\cos\gamma\sin\psi/(r\cos\phi)$ is singular at the
geographic poles $\phi = \pm\pi/2$. This is a coordinate singularity, not a
physical one: $\lambda$ is cyclic, does not enter the corridor constraints
or the glide condition, and can be integrated separately. For any entry
scenario the ground-track extent is bounded by the energy budget, and the
coordinate system can be oriented so the trajectory stays away from the
poles.

The Savageau--Voit polynomial recasting of the thermal subsystem is dropped
entirely. Because the interval pipeline evaluates transcendentals natively,
the density-square-root embedding $y = \sqrt{\rho/\rho_0}$ and the
polynomial surrogate for wall temperature are unnecessary---the raw
physical variables are retained.


%--------------------------------------------------------------------------
\subsection{Frames}\label{sec:frames}
%--------------------------------------------------------------------------

\paragraph{Rotation convention.}
Throughout, $\mat{R}_{\mathcal{A}\mathcal{B}}$ denotes the rotation that
carries coordinates \emph{from} frame~$\mathcal{B}$ \emph{to}
frame~$\mathcal{A}$, i.e.\ $\vvec{v}^{\mathcal{A}} =
\mat{R}_{\mathcal{A}\mathcal{B}}\,\vvec{v}^{\mathcal{B}}$; its columns are the
axes of~$\mathcal{B}$ expressed in~$\mathcal{A}$. Each such matrix is
orthogonal, so $\mat{R}_{\mathcal{A}\mathcal{B}}^{-1} =
\mat{R}_{\mathcal{A}\mathcal{B}}^{\tr} = \mat{R}_{\mathcal{B}\mathcal{A}}$,
and subscripts compose in the usual way,
$\mat{R}_{\mathcal{A}\mathcal{B}}\mat{R}_{\mathcal{B}\mathcal{C}} =
\mat{R}_{\mathcal{A}\mathcal{C}}$.

\begin{itemize}
  \item $\mathcal{E}$ (\textbf{ECEF}): rotating with the constant
    planetary rate $\vvec{\omega}_E = \omega_E\,\hat{z}_{\mathcal{E}}$.

  \item $\mathcal{L}$ (\textbf{LVLH / NED}): at geocentric position
    $(r,\lambda,\phi)$, with basis vectors
    \begin{align*}
      \hat{e}_N &= -\sin\phi\cos\lambda\,\hat{x}_{\mathcal{E}}
                    -\sin\phi\sin\lambda\,\hat{y}_{\mathcal{E}}
                    +\cos\phi\,\hat{z}_{\mathcal{E}},\\
      \hat{e}_E &= -\sin\lambda\,\hat{x}_{\mathcal{E}}
                    +\cos\lambda\,\hat{y}_{\mathcal{E}},\\
      \hat{e}_D &= -\cos\phi\cos\lambda\,\hat{x}_{\mathcal{E}}
                    -\cos\phi\sin\lambda\,\hat{y}_{\mathcal{E}}
                    -\sin\phi\,\hat{z}_{\mathcal{E}}.
    \end{align*}
    The rotation $\mat{R}_{\mathcal{E}\mathcal{L}}$ from $\mathcal{L}$ to
    $\mathcal{E}$ (columns are the above vectors in~$\mathcal{E}$) is a
    function of $(\lambda,\phi)$ alone.

  \item $\mathcal{W}$ (\textbf{Wind frame}): $x$-axis along the
    planet-relative velocity, $z$-axis in the
    $\vvec{V}$--$\hat{e}_D$ plane pointing ``downward'' from the velocity,
    $y$-axis completing the right-handed system. Parameterized by
    flight-path angle~$\gamma$ (positive climbing) and
    heading~$\psi$ (clockwise from North):
    \[
      \hat{V} = \cos\gamma\cos\psi\,\hat{e}_N
              + \cos\gamma\sin\psi\,\hat{e}_E
              - \sin\gamma\,\hat{e}_D.
    \]

  \item $\mathcal{B}$ (\textbf{Body frame}): fixed at the instantaneous
    centre of mass. Related to~$\mathcal{E}$ by the attitude quaternion
    $\vvec{q}$ through the direction-cosine matrix
    $\mat{R}_{\mathcal{E}\mathcal{B}} = \mat{R}(\vvec{q})$, which carries
    body components to~$\mathcal{E}$
    ($\vvec{v}^{\mathcal{E}} = \mat{R}(\vvec{q})\,\vvec{v}^{\mathcal{B}}$):
\end{itemize}
\begin{equation}\label{eq:DCM}
  \mat{R}(\vvec{q})
  = \begin{pmatrix}
      q_0^2+q_1^2-q_2^2-q_3^2 & 2(q_1q_2-q_0q_3) & 2(q_1q_3+q_0q_2)\\
      2(q_1q_2+q_0q_3) & q_0^2-q_1^2+q_2^2-q_3^2 & 2(q_2q_3-q_0q_1)\\
      2(q_1q_3-q_0q_2) & 2(q_2q_3+q_0q_1) & q_0^2-q_1^2-q_2^2+q_3^2
    \end{pmatrix},
\end{equation}
quadratic in~$\vvec{q}$ (its columns are the body axes expressed
in~$\mathcal{E}$).


%--------------------------------------------------------------------------
\subsection{State vector}\label{sec:state-vector}
%--------------------------------------------------------------------------

The full coupled state is
\begin{equation}\label{eq:state-full}
  x = \bigl(\underbrace{r,\lambda,\phi,V,\gamma,\psi}_{6},\;
       \underbrace{q_0,q_1,q_2,q_3,p,q_b,r_b}_{7},\;
       \underbrace{Q_{\mathrm{tot}},T_w,s,m}_{4},\;
       \underbrace{\eta_1,\dots,\eta_{N_m},
                   \dot\eta_1,\dots,\dot\eta_{N_m}}_{2N_m}\bigr),
\end{equation}
with the unit-quaternion constraint $\norm{\vvec{q}}=1$ making the
translational-rotational block effectively 12-dimensional. Total dimension:
$17 + 2N_m$.


%--------------------------------------------------------------------------
\subsection{Translational kinematics}\label{sec:trans-kin}
%--------------------------------------------------------------------------

The planet-relative velocity in the $\mathcal{L}$ (NED) frame is
$\vvec{V}_{\mathcal{L}} =
(V\cos\gamma\cos\psi,\;V\cos\gamma\sin\psi,\;-V\sin\gamma)^\tr$. The
kinematic equations are:
\begin{align}
  \dot r      &= V\sin\gamma, \label{eq:rdot}\\
  \dot\lambda &= \frac{V\cos\gamma\sin\psi}{r\cos\phi},
                                                         \label{eq:lamdot}\\
  \dot\phi    &= \frac{V\cos\gamma\cos\psi}{r}.         \label{eq:phidot}
\end{align}


%--------------------------------------------------------------------------
\subsection{Gravitational acceleration}\label{sec:gravity}
%--------------------------------------------------------------------------

The $J_2$ gravitational potential in geocentric spherical coordinates is
\[
  \Phi_g(r,\phi)
  = -\frac{\mu}{r}\biggl[
      1 - \half J_2\Bigl(\frac{R_e}{r}\Bigr)^{\!2}
          \bigl(3\sin^2\!\phi - 1\bigr)
    \biggr].
\]
Its gradient, resolved in $\mathcal{L}$ (NED), has only radial and
northward components:
\begin{align}
  g_D(r,\phi) &= \frac{\mu}{r^2}\biggl[
    1 + \tfrac{3}{2}\,J_2\Bigl(\frac{R_e}{r}\Bigr)^{\!2}
        \bigl(1-5\sin^2\!\phi\bigr)
  \biggr], \label{eq:gD}\\[4pt]
  g_N(r,\phi) &= -\frac{3\mu}{r^2}\,J_2
    \Bigl(\frac{R_e}{r}\Bigr)^{\!2}
    \sin\phi\cos\phi, \label{eq:gN}\\
  g_E &= 0\quad\text{(by axial symmetry)}.\nonumber
\end{align}

\begin{remark}\label{rem:gravity-spherical}
This reproduces the Cartesian form under the
identifications $r_3/\norm{\vvec{r}} = \sin\phi$ and
$\norm{\vvec{r}} = r$. The spherical form makes the interval evaluation
tighter: $g_D$ and~$g_N$ are explicit functions of $(r,\phi)$, with no
square root of a sum of squares.
\end{remark}


%--------------------------------------------------------------------------
\subsection{Translational dynamics}\label{sec:trans-dyn}
%--------------------------------------------------------------------------

Newton's law in the rotating ECEF frame, resolved along the wind-frame
axes (tangential~$\hat{V}$, normal-in-vertical-plane, lateral), gives
the force equations. Define the aerodynamic force components:
\begin{align*}
  F_V     &= \text{force along } \hat{V},\\
  F_\gamma &= \text{force $\perp\hat{V}$, in the vertical plane},\\
  F_\psi  &= \text{force $\perp$ the vertical plane (lateral)}.
\end{align*}
In the 3-DOF model $F_V=-D$, $F_\gamma=L\cos\sigma$,
$F_\psi=L\sin\sigma$. In the full 6-DOF formulation these are
\emph{derived from the attitude quaternion}---see
\cref{sec:aero-force}.

The translational dynamical equations are:
\begin{align}
  \dot V &= \frac{F_V}{m}
    - g_D\sin\gamma + g_N\cos\gamma\cos\psi
    + \omega_E^2 r\cos\phi
      \bigl[\sin\gamma\cos\phi - \cos\gamma\cos\psi\sin\phi\bigr],
    \label{eq:Vdot}
  \\[6pt]
  \dot\gamma &= \frac{1}{V}\biggl\{
    \frac{F_\gamma}{m}
    - \Bigl(g_D - \frac{V^2}{r}\Bigr)\cos\gamma
    - g_N\sin\gamma\cos\psi
    + 2\omega_E V\cos\phi\sin\psi \nonumber\\
    &\qquad{}+ \omega_E^2 r\cos\phi
      \bigl[\cos\gamma\cos\phi + \sin\gamma\cos\psi\sin\phi\bigr]
    \biggr\}, \label{eq:gamdot}
  \\[6pt]
  \dot\psi &= \frac{1}{V\cos\gamma}\biggl\{
    \frac{F_\psi}{m}
    - \frac{V^2}{r}\cos\gamma\sin\psi\tan\phi
    + g_N\sin\psi \nonumber\\
    &\qquad{}+ 2\omega_E V
      \bigl[\tan\gamma\cos\phi\cos\psi - \sin\phi\bigr]
    + \omega_E^2 r\cos\phi\sin\phi
      \frac{\sin\psi}{\cos\gamma}
    \biggr\}. \label{eq:psidot}
\end{align}

\begin{remark}\label{rem:vinh-reduction}
With $g_N = 0$ (spherical gravity) and
$\omega_E = 0$ (non-rotating planet), these reduce to the classical Vinh
equations for unpowered entry~\cite{vinh1980hypersonic}.
\end{remark}

\begin{remark}\label{rem:psi-singularity}
The singularity $1/\cos\gamma$ in~$\dot\psi$ occurs
at vertical flight ($\gamma = \pm\pi/2$). This is not encountered during
the equilibrium-glide phase and occurs at most instantaneously during a
skip peak. For the interval pipeline, the singular factor is enclosed by
standard interval evaluation of $1/\cos[\gamma]$ provided the interval
$[\gamma]$ does not contain $\pm\pi/2$.
\end{remark}


%--------------------------------------------------------------------------
\subsection{Rotational dynamics}\label{sec:rot-dyn}
%--------------------------------------------------------------------------

The quaternion kinematics, driven by the body rate relative
to~$\mathcal{E}$:
\begin{equation}\label{eq:qdot}
  \dot{\vvec{q}} = \half\,\mat{\Omega}(\vvec{\omega})\,\vvec{q},
\end{equation}
with $\vvec{\omega} = (p,q_b,r_b)^\tr$ the inertial angular velocity in
body-frame components and
\begin{equation}\label{eq:Omega}
  \mat{\Omega}(\vvec{a})
  = \begin{pmatrix}
      0    & -a_1 & -a_2 & -a_3\\
      a_1  &  0   &  a_3 & -a_2\\
      a_2  & -a_3 &  0   &  a_1\\
      a_3  &  a_2 & -a_1 &  0
    \end{pmatrix}.
\end{equation}

Euler's equations, with inertia tensor
$\mat{I} = \mat{I}(s,\eta)$ depending on recession depth and modal
amplitudes:
\begin{equation}\label{eq:euler}
  \mat{I}\,\dot{\vvec{\omega}}
  = \vvec{M} - \vvec{\omega}\times(\mat{I}\,\vvec{\omega})
    - \dot{\mat{I}}\,\vvec{\omega},
\end{equation}
where $\vvec{M} = \vvec{M}_a + \vvec{M}_{\mathrm{rcs}}$ is the total
external moment in~$\mathcal{B}$.


%--------------------------------------------------------------------------
\subsection{Aerodynamic force: from quaternion to
  \texorpdfstring{$(F_V,F_\gamma,F_\psi)$}
  {(FV, Fgamma, Fpsi)}}\label{sec:aero-force}
%--------------------------------------------------------------------------

This section is the bridge between the rotational and translational
dynamics.

\paragraph{Step 1: Body-frame aerodynamic force.}
Given angle of attack~$\alpha$, sideslip~$\beta$, dynamic pressure
$q_{\mathrm{dyn}} = \half\rho V^2$, and reference area~$S$:
\begin{equation}\label{eq:Fa-body}
  \vvec{F}_a^{\mathcal{B}}
  = q_{\mathrm{dyn}}\,S\,
    \bigl(-C_A(\alpha,\beta,M),\;
           C_Y(\alpha,\beta,M),\;
          -C_N(\alpha,\beta,M)\bigr)^\tr.
\end{equation}

\paragraph{Step 2: Transform to LVLH.}
Define the body-to-LVLH rotation
\begin{equation}\label{eq:RLB}
  \mat{R}_{\mathcal{L}\mathcal{B}}({\vvec{q}},\lambda,\phi)
  \;:=\;
  \mat{R}_{\mathcal{E}\mathcal{L}}^\tr(\lambda,\phi)\;
  \mat{R}(\vvec{q})
  \;=\;
  \mat{R}_{\mathcal{L}\mathcal{E}}\,\mat{R}_{\mathcal{E}\mathcal{B}},
\end{equation}
so that $\vvec{v}^{\mathcal{L}} =
\mat{R}_{\mathcal{L}\mathcal{B}}\,\vvec{v}^{\mathcal{B}}$. Then the
aerodynamic force in~$\mathcal{L}$ is
\begin{equation}\label{eq:Fa-LVLH}
  \vvec{F}_a^{\mathcal{L}}
  = \mat{R}_{\mathcal{L}\mathcal{B}}\,
    \vvec{F}_a^{\mathcal{B}}.
\end{equation}

\paragraph{Step 3: Project onto wind axes.}
The wind-axis triad, expressed in~$\mathcal{L}$, is
\begin{align}
  \hat{V}_{\mathcal{L}} &= (\cos\gamma\cos\psi,\;
                      \cos\gamma\sin\psi,\;-\sin\gamma)^\tr,
                                                    \label{eq:Vhat}\\
  \hat{n}_\gamma &= (-\sin\gamma\cos\psi,\;
                      -\sin\gamma\sin\psi,\;
                      -\cos\gamma)^\tr, \label{eq:ngam}\\
  \hat{n}_\psi   &= \hat{V}_{\mathcal{L}} \times \hat{n}_\gamma.
                                                    \label{eq:npsi}
\end{align}
Then:
\begin{equation}\label{eq:F-projections}
  F_V = \vvec{F}_a^{\mathcal{L}}\cdot\hat{V}_{\mathcal{L}},\qquad
  F_\gamma = \vvec{F}_a^{\mathcal{L}}\cdot\hat{n}_\gamma,\qquad
  F_\psi = \vvec{F}_a^{\mathcal{L}}\cdot\hat{n}_\psi.
\end{equation}

\paragraph{Step 4: Identification with lift, drag, bank angle.}
In the point-mass limit where the attitude is in trim,
$F_V = -D$, $F_\gamma = L\cos\sigma$, $F_\psi = L\sin\sigma$.
In the full 6-DOF formulation $D$, $L$, $\sigma$ are \emph{not independent
inputs}: they are determined by the attitude quaternion~$\vvec{q}$, the
aerodynamic coefficients, and the velocity direction $(\gamma,\psi)$
through Eqs.~(\ref{eq:Fa-body})--(\ref{eq:F-projections}).

\begin{remark}\label{rem:bank-angle}
The effective bank angle can be extracted as
$\sigma = \operatorname{atan2}(F_\psi, F_\gamma)$,
an algebraic function of
$({\vvec{q}},\lambda,\phi,\gamma,\psi,\alpha,\beta)$. Under interval
arithmetic, $\sigma$ is evaluated as an interval $\operatorname{atan2}$, a
standard operation in validated numerics~\cite{moore2009introduction}.
\end{remark}

\begin{remark}\label{rem:aoa-sideslip}
The angle of attack~$\alpha$ and sideslip~$\beta$ are
determined by the quaternion and velocity direction. Let
$\vvec{V}^{\mathcal{B}} = \mat{R}_{\mathcal{L}\mathcal{B}}^{\tr}\,
\vvec{V}_{\mathcal{L}}$ be the velocity in body-frame components. Then
$\alpha = \operatorname{atan2}(-V_z^{\mathcal{B}}, -V_x^{\mathcal{B}})$
and $\beta = \arcsin(V_y^{\mathcal{B}}/V)$.
\end{remark}


%--------------------------------------------------------------------------
\subsection{Equilibrium-glide slow manifold}\label{sec:glide-manifold}
%--------------------------------------------------------------------------

The equilibrium-glide condition is the slow manifold of the translational
dynamics under the thin-atmosphere singular perturbation
$\eps_{\mathrm{atm}} = H_s/R_e \to 0$.

Setting $\dot\gamma = 0$ and neglecting Earth rotation ($\omega_E = 0$) in
Eq.~(\ref{eq:gamdot}):
\[
  \frac{F_\gamma}{m}
  = \Bigl(g_D - \frac{V^2}{r}\Bigr)\cos\gamma
    + g_N\sin\gamma\cos\psi.
\]
For small flight-path angle ($\gamma\ll 1$) and spherical gravity
($g_N = 0$, $g_D = \mu/r^2$):
\begin{equation}\label{eq:glide-3dof}
  \frac{F_\gamma}{m} = g - \frac{V^2}{r},
  \qquad\text{i.e.}\quad
  \frac{L\cos\sigma}{m(g - V^2/r)} = 1
  \quad(\Lambda = 1).
\end{equation}

In the full 6-DOF formulation, $F_\gamma$ is determined by the quaternion
via Eqs.~(\ref{eq:Fa-body})--(\ref{eq:F-projections}). The slow-manifold
condition therefore constrains the quaternion directly:
\begin{multline}\label{eq:glide-full}
  \bigl[\mat{R}_{\mathcal{L}\mathcal{B}}\,
         \vvec{F}_a^{\mathcal{B}}\bigr]
  \cdot\hat{n}_\gamma \\
  = m\Bigl(g_D(r,\phi) - \frac{V^2}{r}\Bigr)\cos\gamma
    + m\,g_N(r,\phi)\sin\gamma\cos\psi,
\end{multline}
a single scalar constraint on the coupled translational and rotational state
(dimension~$13$); it defines the codimension-1 slow manifold~$\Lambda$.

\begin{remark}\label{rem:manifold-attitude}
The slow manifold is \emph{not} a constraint on
$\sigma$ alone---it constrains a specific component of
$\mat{R}_{\mathcal{L}\mathcal{B}} \vvec{F}_a^{\mathcal{B}}$, which
depends on the full quaternion and on the aerodynamic angles
$(\alpha,\beta)$ implicitly. The ``bank angle'' that appears in the
point-mass glide condition is a derived quantity, and the slow manifold
constrains the attitude directly.
\end{remark}

\begin{remark}\label{rem:j2-glide}
Equation~(\ref{eq:glide-full}) retains the $J_2$
contributions through $g_D(r,\phi)$ and $g_N(r,\phi)$. The simplified
form $\Lambda = L\cos\sigma/[m(g-V^2/r)] = 1$ is the leading-order
reduction under spherical gravity ($J_2 = 0$) and small flight-path angle.
The $J_2$ correction to the glide balance is $O(J_2)$, ranging from
${\approx}\,0.16\%$ of the leading gravitational term near the equator to
${\approx}\,0.65\%$ near the poles (through the $(1-5\sin^2\!\phi)$ factor);
for the interval pipeline it is enclosed exactly.
\end{remark}

\begin{remark}\label{rem:j2-skip}
The $J_2$ contribution is especially significant
during skip maneuvers. When the vehicle leaves the slow manifold (the
defining condition of a boundary-layer bubble), three $J_2$ effects
compound:
\begin{enumerate}
  \item The northward gravitational component $g_N(r,\phi)$ enters the
    heading rate $\dot\psi$ (Eq.~\ref{eq:psidot}), producing a
    latitude-dependent heading drift during each skip that is absent under
    spherical gravity. Northward and southward skips are not symmetric.
  \item The radial component $g_D(r,\phi)$ varies with latitude by up to
    $\approx 0.65\%$ of the leading term through the $(1-5\sin^2\!\phi)$
    factor, modifying the
    effective centrifugal balance $g_D - V^2/r$ that governs the
    re-entry angle after each skip peak.
  \item Over $K$ skips the ground track may span a significant latitude
    range, so the $J_2$ corrections at successive re-entries are evaluated
    at different~$\phi$, and the effects are cumulative.
\end{enumerate}
For the interval pipeline, these effects are enclosed automatically: the
interval evaluations of $g_D([r],[\phi])$ and $g_N([r],[\phi])$ widen with
the latitude interval, which is itself widest during a skip (where the
ground track sweeps the most latitude). Retaining $J_2$ in the dynamics
is therefore not a refinement but a necessity for tight certified bounds on
skip trajectories.
\end{remark}


%--------------------------------------------------------------------------
\subsection{Thermal subsystem (native form)}\label{sec:thermal}
%--------------------------------------------------------------------------

The atmosphere density profile is $\rho(h)$ from the NRLMSIS~2.0
model~\cite{emmert2021nrlmsis}, with $h = r - R_e$. Under interval arithmetic,
$\rho(h)$ is enclosed by evaluating the model's analytic approximation (or
a validated piecewise-exponential envelope) with interval
argument~$[h]$.

Stagnation-point heating (Sutton--Graves~\cite{sutton1971general}, within its
stated validity envelope: axisymmetric blunt bodies, chemical equilibrium,
enthalpy 2.3--116.2\,MJ/kg, pressure 0.001--100\,atm):
\begin{equation}\label{eq:Qdot}
  \dot Q_{\mathrm{tot}} = k_Q\sqrt{\rho/R_n}\,V^3.
\end{equation}
Wall temperature:
\begin{equation}\label{eq:Twdot}
  \dot T_w
  = \frac{\dot Q_{\mathrm{tot}} - (T_w - T_\infty)/\tau_{\mathrm{cond}}}
         {C_w}.
\end{equation}
No embedding, no recasting. The transcendentals $\sqrt{\rho}$ and $V^3$
are evaluated directly as interval operations.

\begin{remark}\label{rem:thermal-embedding}
The density-square-root embedding
$y = \sqrt{\rho/\rho_0}$ of the original formulation was introduced to
render the heating correlation polynomial in the states (the
Savageau--Voit recasting strategy~\cite{savageau1987recasting}). Since the interval
pipeline evaluates $\sqrt{\cdot}$ natively, this auxiliary state is
eliminated, and with it the associated wrapping from carrying an extra
dimension.
\end{remark}


%--------------------------------------------------------------------------
\subsection{Ablation}\label{sec:ablation}
%--------------------------------------------------------------------------

Recession rate coupled to heat flux; mass loss from char-layer blowoff:
\begin{equation}\label{eq:ablation}
  \dot s = \alpha_s\,\dot Q_{\mathrm{tot}},\qquad
  \dot m = -\dot m_{\mathrm{abl}}.
\end{equation}


%--------------------------------------------------------------------------
\subsection{Modal / structural elasticity}\label{sec:modal}
%--------------------------------------------------------------------------

Per mode~$i$, with piston-theory forcing and temperature-dependent natural
frequency:
\begin{equation}\label{eq:modal}
  \ddot\eta_i + 2\zeta_i\,\omega_i(T_w)\,\dot\eta_i
  + \omega_i(T_w)^2\,\eta_i
  = Q_i(x,\eta,\dot\eta).
\end{equation}
Piston theory is a double expansion in thickness ratio and inverse Mach
number valid when the contributions of order $\delta/M^3$ and
$\delta^2/M^2$ are negligible~\cite{ashley1956piston,lighthill1953oscillating}; its
applicability conditions must be checked at each atmospheric pass. At
skip peaks with high angle of attack and strong shock interaction, piston
theory may fail, and the modal forcing should be flagged as uncertain
(widened interval) in those regimes.


%--------------------------------------------------------------------------
\subsection{Density model and interval envelope}\label{sec:density}
%--------------------------------------------------------------------------

\emph{To be completed:} functional form of the atmospheric density profile
suitable for interval evaluation. The NRLMSIS~2.0 model provides the
reference; the interval pipeline requires either a closed-form analytic
approximation with certified error bounds, or a piecewise-exponential
envelope
\[
  \rho_{\mathrm{lower}}(h) \le \rho(h) \le \rho_{\mathrm{upper}}(h)
\]
valid across the altitude range of the corridor. The width
$\rho_{\mathrm{upper}} - \rho_{\mathrm{lower}}$ at each altitude is the
primary driver of the aerodynamic interval widths.


%--------------------------------------------------------------------------
\subsection{Summary of coordinate choices}\label{sec:summary-table}
%--------------------------------------------------------------------------

\begin{center}
\begin{tabular}{@{}lll@{}}
\toprule
\textbf{Feature} &
\textbf{Spherical-LVLH + Quat.} &
\textbf{Cartesian ECEF + Quat.} \\
\midrule
Corridor constraints &
  Box in $(r,V)$ &
  Nonlinear: $h = \norm{\vvec{r}} - R_e$ \\
Glide condition &
  Native: Eq.~(\ref{eq:glide-3dof})/(\ref{eq:glide-full}) &
  Derived from $(u,v,w)$, $\norm{\vvec{r}}$ \\
Position singularity &
  Pole in $\dot\lambda$ &
  None \\
Dynamics &
  Transcendental ($\sin$, $\cos$) &
  Polynomial (quadratic in $\vvec{q}$) \\
Wrapping source &
  Trig dependency &
  Corridor intersection \\
Density evaluation &
  $\rho(r-R_e)$: direct &
  $\rho(\norm{\vvec{r}}-R_e)$: extra $\sqrt{\cdot}$ \\
\bottomrule
\end{tabular}
\end{center}

\smallskip\noindent
The spherical-LVLH choice concentrates wrapping in the dynamics
(trigonometric evaluations), where it is distributed across every time step
and can be controlled by step-size refinement. The Cartesian choice
concentrates wrapping in the corridor intersection, where it is applied
once per intersection and inflates the entire reachable-set bound. For the
profile decomposition---which is organized around the corridor and the
glide manifold---the spherical choice produces tighter certified bounds.

% =====================================================================
%  Thesis proposal, Appendix: Concentration along the time axis
%  Owns: apx:time
% =====================================================================

\section{Concentration along the time axis}
\label{apx:time}

This appendix supports \cref{sec:profile}. It is a sketch, and it says
where it stops. Two elementary bounds and the compactness they give are
proved; the application of Lions' lemma is a statement about a sequence of
densities and is exact; the identification of its alternatives with the
flight regimes, and the passage from the dissipation back to the state,
are stated as what remains to be shown, and gathered in
\cref{sec:time-open}.

%--------------------------------------------------------------------------
\subsection{Setting, and two elementary bounds}\label{sec:time-setting}
%--------------------------------------------------------------------------

Throughout, $\eps = \eps_{\mathrm{atm}}$. The corridor $\calX$ is a compact
set of states on which $V \ge V_{\min} > 0$,
$\abs{\gamma}\le\gamma_{\max}<\pi/2$ and $\abs{\phi}\le\phi_{\max}<\pi/2$.
A \emph{trajectory} is an absolutely continuous
$x_\eps\colon[0,T]\to\calX$ whose translational part
satisfies~\crefrange{eq:rdot}{eq:phidot} and~\crefrange{eq:Vdot}{eq:psidot}
for almost every~$t$, with the specific aerodynamic force
$(F_V,F_\gamma,F_\psi)/m$ selected measurably from the enclosure. Nothing is
assumed about how the force arises: the bank history, the attitude and the
selection from the interval field are all free. One assumption is made on
the enclosure.

\begin{assum}[Dissipative enclosure]\label{asm:dissipative}
There is a constant $E_{\max}$ such that every selection satisfies
$F_V \le 0$ and $(F_\gamma^2 + F_\psi^2)^{1/2} \le E_{\max}\abs{F_V}$.
\end{assum}

The first inequality says that the aerodynamic force never does positive
work on the vehicle, the second that lift is bounded by a multiple of drag.
For a lift-and-drag enclosure they read $[C_D]\subset(0,\infty)$ and
$[C_L]/[C_D]\subset[-E_{\max},E_{\max}]$. A propulsive mode adds a bounded
source to what follows and is not considered here.

\begin{lem}[Energy balance]\label{lem:energy}
Let
\[
  \mathcal{E} \;=\; \tfrac12 V^2 + \Phi_g(r,\phi)
                    - \tfrac12\,\omega_E^2 r^2\cos^2\!\phi
\]
be the specific mechanical energy in the rotating frame, with $\Phi_g$ the
potential of \cref{sec:gravity}. Along every trajectory
\begin{equation}\label{eq:dissipation}
  \dot{\mathcal{E}} \;=\; \frac{F_V}{m}\,V \;=:\; -P \;\le\; 0 ,
\end{equation}
so that $\int_0^T P\dd t = \mathcal{E}(x(0)) - \mathcal{E}(x(T))
\le \Delta\mathcal{E}$, where
$\Delta\mathcal{E} = \max_{\calX}\mathcal{E} - \min_{\calX}\mathcal{E}$.
\end{lem}

\begin{proof}
Differentiate along \crefrange{eq:rdot}{eq:phidot} and \cref{eq:Vdot}. The
terms in $g_D$ and $g_N$ cancel against
$\partial_r\Phi_g\,\dot r + \partial_\phi\Phi_g\,\dot\phi$, by the
definition of $g_D$ and $g_N$ as the gradient of $\Phi_g$, and the
centrifugal term of \cref{eq:Vdot} against the derivative of
$\tfrac12\omega_E^2r^2\cos^2\!\phi$. The Coriolis force and the two normal
components of the aerodynamic force are perpendicular to the velocity and
do not appear in \cref{eq:Vdot} at all. The identity is checked by computer
algebra with oblateness and rotation retained.
\end{proof}

\begin{lem}[Impulse bound]\label{lem:impulse}
Under \cref{asm:dissipative} the aerodynamic acceleration
$\vvec{a} = \vvec{F}_a/m$ satisfies $\abs{\vvec{a}} \le C_a P$ with
$C_a = \sqrt{1+E_{\max}^2}\,/\,V_{\min}$. Hence for every interval
$I\subseteq[0,T]$,
\[
  \int_I \abs{\vvec{a}}\dd t \;\le\; C_a\int_I P\dd t .
\]
\end{lem}

\begin{proof}
$\abs{\vvec{a}}^2 = (F_V^2 + F_\gamma^2 + F_\psi^2)/m^2
\le (1+E_{\max}^2)\,F_V^2/m^2$ and $\abs{F_V}/m = P/V \le P/V_{\min}$.
\end{proof}

The second lemma is what makes the dissipation the right quantity to follow.
Lift cannot be had without drag, so wherever no energy is lost the
aerodynamics does nothing, and the motion is that of the gravity field in
the rotating frame. Whatever a pass does to the state, the dissipation sees
it.

%--------------------------------------------------------------------------
\subsection{Compactness in slow time}\label{sec:time-slow}
%--------------------------------------------------------------------------

Let $\eps_n\to 0$ and let $x_n$ be a trajectory for each~$n$. Write
$\mathcal{E}_n(t) = \mathcal{E}(x_n(t))$; $\pi_n = P_n\dd t$, a positive
measure on $[0,T]$ of mass at most $\Delta\mathcal{E}$; and
$\vvec{y}_n(t)$, $\vvec{v}_n(t)\in\R^3$ for the position and the
planet-relative velocity in planet-fixed Cartesian components, in which the
chart singularities play no part.

\begin{prop}[Skeleton and impulses as a limit]\label{prop:slow}
Under \cref{asm:dissipative} there is a subsequence along which:
\begin{enumerate}[label=(\roman*),leftmargin=*,itemsep=2pt]
  \item $\mathcal{E}_n(t)\to\mathcal{E}_0(t)$ for every $t\in[0,T]$, with
    $\mathcal{E}_0$ non-increasing, and $\pi_n\rightharpoonup\pi_0$
    weakly-$*$, where $\pi_0$ is the measure with distribution function
    $\mathcal{E}_0(0)-\mathcal{E}_0$;
  \item $\pi_0 = p\dd t + \sum_j a_j\,\delta_{t_j} + \pi_c$, with
    $p\in L^1(0,T)$, at most countably many atoms $a_j>0$, and $\pi_c$
    singular and without atoms; and for every $\delta>0$
    \[
      \#\{\,j : a_j \ge \delta\,\} \;\le\; \Delta\mathcal{E}/\delta ;
    \]
  \item on every closed interval on which $\mathcal{E}_0$ is continuous,
    the convergence $\mathcal{E}_n\to\mathcal{E}_0$ is uniform;
  \item $\vvec{y}_n\to\vvec{y}_0$ uniformly, $\vvec{v}_n(t)\to\vvec{v}_0(t)$
    for every~$t$, $\dot{\vvec{y}}_0 = \vvec{v}_0$ almost everywhere, and
    for $0\le s<t\le T$
    \begin{equation}\label{eq:bv}
      \abs{\vvec{v}_0(t) - \vvec{v}_0(s)}
      \;\le\; C_g\,(t-s) + C_a\,\pi_0([s,t]) ,
    \end{equation}
    with $C_g$ a bound on the gravitational, Coriolis and centrifugal
    accelerations over the corridor.
\end{enumerate}
\end{prop}

\begin{proof}
(i) The functions $\mathcal{E}_n$ are non-increasing by \cref{lem:energy}
and uniformly bounded on the compact corridor, so Helly's selection theorem
gives a subsequence converging at every point; the measures $\pi_n$ are
their Stieltjes measures, and pointwise convergence of the distribution
functions at the continuity points of the limit, with bounded total mass,
is weak-$*$ convergence. (ii) is the Lebesgue decomposition of $\pi_0$, and
the count follows from $\sum_j a_j\le\pi_0([0,T])\le\Delta\mathcal{E}$.
(iii) is P\'olya's theorem: monotone functions converging pointwise to a
continuous limit on a compact interval converge uniformly. (iv) In
planet-fixed components Newton's law reads
$\dot{\vvec{v}}_n = \vvec{a}_n + \vvec{g}(\vvec{y}_n)
 - 2\,\vvec{\omega}_E\times\vvec{v}_n
 - \vvec{\omega}_E\times(\vvec{\omega}_E\times\vvec{y}_n)$,
so by \cref{lem:impulse}
\[
  \abs{\vvec{v}_n(t)-\vvec{v}_n(s)}
  \;\le\; C_g\,(t-s) + C_a\,\pi_n([s,t]) .
\]
The velocities therefore have total variation at most
$C_g T + C_a\,\Delta\mathcal{E}$, uniformly in~$n$, and Helly's theorem for
functions of bounded variation gives pointwise convergence along a further
subsequence. The positions are uniformly Lipschitz and converge uniformly
by Arzel\`a--Ascoli, and $\dot{\vvec{y}}_0=\vvec{v}_0$ by dominated
convergence. \Cref{eq:bv} follows on letting $n\to\infty$, since
$\limsup_n\pi_n(F)\le\pi_0(F)$ for closed~$F$.
\end{proof}

Three things can be read off \cref{eq:bv}. The limiting velocity has bounded
variation, and it can jump only at an atom of $\pi_0$, by at most
$C_a a_j$: an atom is an impulse, and its size bounds the impulse. On an
interval that $\pi_0$ does not charge, $\int\abs{\vvec{a}_n}\dd t\to 0$ and
the limit is an arc of the gravity field alone. And where $\pi_0$ has a
density~$p$ the aerodynamic forcing survives as an acceleration of
magnitude at most $C_a\,p$, which is the glide. This is the
skeleton-and-impulses picture as a statement about
limits, valid for every control and every selection, and it is as far as
compactness alone goes. It does not exclude the part $\pi_c$, and it says
nothing yet about what happens \emph{inside} an atom.

%--------------------------------------------------------------------------
\subsection{The fast-time axis and Lions' lemma}\label{sec:time-fast}
%--------------------------------------------------------------------------

To see inside an atom the time axis is dilated. Fix a scale $h_n\to 0$---in
the body, $h_n=\sqrt{\eps_n}$---suppose the total dissipation stays bounded
below, $\pi_n([0,T])\ge m_0>0$, and set
\[
  \nu_n(\tau) \;=\; \frac{h_n\,P_n(h_n\tau)}{\pi_n([0,T])}
  \qquad(\tau\in\R),
\]
extended by zero outside $[0,T/h_n]$. Then $\nu_n\ge 0$ and
$\int_\R\nu_n\dd\tau = 1$, and the concentration-compactness
lemma~\cite[Lemma~I.1]{lions1984concentration1} applies with no further
hypothesis. There is a subsequence along which one of the following holds.
\begin{description}[leftmargin=1.5em,itemsep=2pt]
  \item[Compactness.] There are $\tau_n\in\R$ such that for every
    $\eta>0$ there is $R<\infty$ with
    $\int_{\tau_n-R}^{\tau_n+R}\nu_n\dd\tau \ge 1-\eta$ for all~$n$.
  \item[Vanishing.] For every $R<\infty$,
    $\;\sup_{y\in\R}\int_{y-R}^{y+R}\nu_n\dd\tau \to 0$.
  \item[Dichotomy.] There is $\alpha\in(0,1)$ such that for every $\eta>0$
    there are nonnegative $\nu_n^1,\nu_n^2$ with
    $\norm{\nu_n - \nu_n^1 - \nu_n^2}_{L^1}\le\eta$,
    $\abs{\int\nu_n^1 - \alpha}\le\eta$,
    $\abs{\int\nu_n^2 - (1-\alpha)}\le\eta$, and
    $\operatorname{dist}(\operatorname{supp}\nu_n^1,
    \operatorname{supp}\nu_n^2)\to\infty$.
\end{description}

Each has a reading. In the first, all but a fraction $\eta$ of the
dissipation occurs within a time $R\,h_n$ of $t_n = h_n\tau_n$: a single
pass, and $\pi_0$ is one atom. In the second no window of width $R\,h_n$
holds a fixed share. This is the case whenever the drag power is bounded
uniformly in~$n$, since then a window holds at most $2RC\,h_n/m_0$; the
limit $\pi_0$ then has a bounded density and no atom. In the third the
dissipation splits into pieces that separate on the fast scale: two passes
a positive time apart, or a pass followed by a glide. Applying the lemma
again to each piece, as in the construction of a profile decomposition,
extracts the concentrated pieces one at a time; for a threshold~$\delta$
the process stops after at most $\Delta\mathcal{E}/\delta$ steps, by
\cref{prop:slow}(ii), and leaves a remainder in which no window holds more
than about~$\delta$. The remainder is small in the sense of the second
alternative, in the norm $\sup_y\int_{y-R}^{y+R}\abs{\,\cdot\,}\dd\tau$,
and not in $L^1$: on a glide it is most of the dissipation.

\begin{remark}[Which alternative occurs is decided by the corridor]
\label{rem:regimes}
Whether a family concentrates at all is a statement about how its corridor
scales with $\eps$, and belongs among the hypotheses. If the dynamic
pressure and load limits are held fixed as $\eps\to 0$, the drag power is
bounded, every sequence vanishes, and the limit has no atoms: the
trajectories are glide and phugoid, and the fast-scale structure is
\emph{oscillation}, for which the natural tool is averaging over the
invariant measures of the fast flow~\cite{artstein1999invariant}.
Concentration requires loads that grow as $\eps\to 0$, as in the classical
matched-asymptotic analysis of skip entry at a fixed entry
angle~\cite{shi1971skip}, where the layer has thickness $O(\eps)$ in
altitude. The physical value $\eps\approx 1.1\times10^{-3}$, with entry
angles of a few degrees, lies between the two: by \cref{eq:turning} a pass
that turns the flight path through $8^\circ$ at a vertical lift-to-drag
ratio of two dissipates about an eighth of the kinetic energy, and a
computed pass lasts about a fifth of the slow time.
The threshold in \cref{eq:kmax} is what makes the decomposition well posed
at fixed $\eps$ without deciding the regime. Passes that dissipate more
than $\delta$ are bubbles, and there are at most $\Delta\mathcal{E}/\delta$
of them whatever $\eps$ is; everything else is residual, and has to be
bounded as oscillation.
\end{remark}

\begin{remark}[The scale has to be chosen for each piece]
Compactness at the scale $h_n$ forces a single atom of $\pi_0$, but
vanishing at that scale does not exclude one: dissipation that concentrates
at a coarser scale $h_n'\gg h_n$, $h_n'\to 0$, vanishes at the scale $h_n$
and still produces an atom. Conversely a pass shorter than $h_n$ is compact
at that scale and its profile is a Dirac mass. The lemma classifies what
happens at one scale; the width of each piece has to be found from the
piece. That is the content of the limit-case form of the
principle~\cite{lions1985limit1}, in which compactness is recovered up to a
translation \emph{and} a dilation.
\end{remark}

%--------------------------------------------------------------------------
\subsection{Profiles and their orthogonality}\label{sec:time-profiles}
%--------------------------------------------------------------------------

For an atom $(t_j,a_j)$ of $\pi_0$ choose centers $t_n^j\to t_j$ and widths
$h_n^j\to 0$ such that the windows
$[t_n^j - R\,h_n^j,\, t_n^j + R\,h_n^j]$ capture $a_j$ as $R\to\infty$,
uniformly in~$n$, with $h_n^j$ the smallest scale at which this holds. The
rescaled densities $s\mapsto h_n^j\,P_n(t_n^j + h_n^j s)$ are then tight on
the line and converge weakly-$*$ along a subsequence to a measure $p^j$ of
mass $a_j$, the \emph{profile} of the pass. Two profiles $j\ne k$ are
\emph{orthogonal} in the sense of~\cite[(1.18)]{bahouri1999high} when
\[
  \Bigl|\log\frac{h_n^j}{h_n^k}\Bigr| \to \infty
  \qquad\text{or}\qquad
  h_n^j = h_n^k \ \text{ and } \
  \frac{\abs{t_n^j - t_n^k}}{h_n^j}\to\infty .
\]
Atoms at distinct times are orthogonal by the second alternative, with
nothing to prove. For orthogonal profiles the dissipation is additive,
\begin{equation}\label{eq:additivity}
  \pi_n([0,T])
  \;=\; \sum_{j=1}^{J} a_j \;+\; r_n^J([0,T]) \;+\; o(1)
  \qquad(n\to\infty),
\end{equation}
with $r_n^J$ the dissipation outside the $J$ windows. This is the analogue
of the energy identity~\cite[(1.25)]{bahouri1999high}, and it is immediate,
because the dissipation is a measure and the windows are eventually
disjoint. The analogue of the nonlinear
statement~\cite[(1.26)]{bahouri1999high}---that the solution is, up to a
small error, the superposition of the profiles' separate evolutions---is
not immediate, and here it is not a superposition at all. An ordinary
differential equation composes: the state leaving pass~$j$ is the datum of
the arc that leads to pass~$j+1$. What replaces (1.26) is a factorization
of the map from entry data to exit data through the individual pass maps
and the outer flow, with a remainder that is explicit and uniform over the
selections. That is objective~4.

%--------------------------------------------------------------------------
\subsection{Joint profiles where scales meet}\label{sec:time-joint}
%--------------------------------------------------------------------------

Within a single pass the attitude and elastic motions are excited at their
own scales, $h^{\mathrm{att}} = \eps_{\mathrm{att}}$ and
$h^{\mathrm{ela}} = \eps_{\mathrm{ela}}$ in units of the slow time. They
share the center of the pass, so the second alternative of (1.18) is not
available and orthogonality has to come from the first: the ratio of the
scales must diverge along the family. Where it does, the faster motion
averages out over the slower. Where two periods remain comparable---where
two bands cross inside the corridor---the corresponding profiles cannot be
separated, and the decomposition carries one joint profile on the coupled
subsystem. The check is pairwise and depends on the vehicle. For generic
parametrizations the first structural period lies between
$6\,\mathrm{ms}$ and $0.18\,\mathrm{s}$ and a pass lasts about
$150\,\mathrm{s}$; with the attitude period taken between $0.1$ and
$10\,\mathrm{s}$, the atmospheric scale is separated from the other two by
more than a decade, and the attitude and elastic scales meet only for a
very slender, flexible body with a fast attitude loop. These are orders of
magnitude, to be redone for any vehicle to which the result is applied.

%--------------------------------------------------------------------------
\subsection{What is not proved}\label{sec:time-open}
%--------------------------------------------------------------------------

\begin{enumerate}[leftmargin=*,itemsep=3pt]
  \item \emph{The singular continuous part.} \Cref{prop:slow} allows
    $\pi_c\ne 0$: dissipation carried by a set of times of measure zero
    without atoms. Countably many passes that accumulate at a finite time,
    with summable dissipation---a sequence of decaying skips that ends in a
    glide, as the bounces of a ball end in rolling---are atoms and are
    already covered. A genuinely diffuse singular part would have to be
    excluded by a lower bound on the separation of passes in terms of the
    energy each dissipates. This step is open.
  \item \emph{Vanishing is glide.} That the part of $\pi_0$ with a density
    corresponds to motion near the glide balance, and how near, is the
    estimate of \cref{apx:residual}. It is an averaging estimate for a
    weakly damped oscillator, not a boundary-layer estimate, and it is not
    yet uniform over the selections.
  \item \emph{Compactness of the profiles over the family.} The profile
    $p^j$ and the impulse set it generates depend on the selection.
    \Cref{lem:impulse} bounds the diameter of the impulse set by $C_a a_j$;
    a description of the set itself is the inner problem of the classical
    theory, to be solved over the interval field.
  \item \emph{Composition.} The factorization that replaces (1.26), with
    its remainder, at each of the three scales and under interval
    disturbances.
\end{enumerate}

% =====================================================================
%  Thesis proposal, Appendix: Skeleton residual and bubble budget
%  Owns: apx:residual
% =====================================================================

\section{Skeleton residual and bubble budget}
\label{apx:residual}

This appendix supports \cref{sec:bound}. It gives the two constants behind
\cref{eq:target} that presently exist, the lag behind the glide and the
count of passes, together with the equation that connects glide, phugoid
and skip; it says what has been computed against them, and states the range
in which each holds. The
identities are checked by computer algebra, and the leading-order
statements as limits of exact expressions, against the same definition of
the field that is evaluated numerically; each comparison with computation
was run against a failure criterion stated beforehand.

%--------------------------------------------------------------------------
\subsection{The planar glide}\label{sec:residual-glide}
%--------------------------------------------------------------------------

For spherical gravity and no rotation the altitude, speed and flight-path
angle decouple from the rest of~\crefrange{eq:rdot}{eq:psidot}:
\[
  \dot r = V\sin\gamma,\qquad
  \dot V = -\frac{D}{m} - g\sin\gamma,\qquad
  V\dot\gamma = \frac{L\cos\sigma}{m} - \Bigl(g - \frac{V^2}{r}\Bigr)\cos\gamma ,
\]
with $D/m = \rho V^2/(2\beta)$ for a ballistic coefficient~$\beta$, an
exponential density of scale height~$H_s$, and $g = \mu/r^2$. Two
abbreviations recur: the vertical lift-to-drag ratio and the net
gravitational acceleration,
\[
  E_v = \frac{L}{D}\cos\sigma, \qquad G = g - \frac{V^2}{r} .
\]
On the equilibrium glide $\gamma = 0$ and $L\cos\sigma/m = G$, so
$D/m = G/E_v$. As the speed bleeds off the glide altitude descends, and to
leading order in $\eps = H_s/R_e$ its rate is
\begin{equation}\label{eq:glide-descent}
  \dot h_{\mathrm{g}} \;=\; -\,\frac{2H_s\,g}{V E_v} .
\end{equation}

%--------------------------------------------------------------------------
\subsection{The phugoid}\label{sec:residual-phugoid}
%--------------------------------------------------------------------------

Linearized about a point of the glide with the coefficients frozen, the
three equations have the characteristic polynomial
$\lambda^3 + a_2\lambda^2 + a_1\lambda + a_0$, in which
$a_2 = 2G/(E_vV)$ exactly, $H_s\,a_1\to G$ and $a_0/a_1\to -2V/(rE_v)$ as
$H_s\to 0$. Its roots are one real root $\lambda_{\mathrm{s}}$ and a
complex pair $-\varsigma\pm i\omega$, the \emph{phugoid}: the oscillation of
altitude against flight-path angle about the glide. The sum of the roots is
$-a_2$, so $\varsigma = (a_2 + \lambda_{\mathrm{s}})/2$ exactly, and to
leading order $\lambda_{\mathrm{s}} = -a_0/a_1 = 2V/(rE_v)$. Hence
\begin{equation}\label{eq:phugoid}
  \omega^2 \;=\; \frac{G}{H_s}, \qquad
  \varsigma \;=\; \frac{g}{V E_v},
\end{equation}
each up to a relative correction of order~$\eps$. The real root is positive
and is not an instability. It is the rate at which a frozen linearization
registers that the glide itself is descending.

The two quantities in \cref{eq:phugoid} scale differently, and everything
in this appendix follows from that. The frequency grows like
$\eps^{-1/2}$. The damping rate does not contain $H_s$ at all; by
\cref{eq:glide-descent} it is $\abs{\dot h_{\mathrm{g}}}/(2H_s)$, half the
rate at which the glide descends through scale heights. The damping ratio is
therefore
\begin{equation}\label{eq:zeta}
  \zeta \;=\; \frac{\varsigma}{\omega}
  \;=\; \frac{\sqrt{g R_e}}{V E_v\sqrt{1 - V^2/(gr)}}\;\sqrt{\eps} ,
\end{equation}
of order $\sqrt{\eps}$. For $L/D = 2$ at $6$\,km/s and the terrestrial
scale height the period is about $265$\,s, the decay time
$1/\varsigma\approx 1250$\,s and $\zeta\approx 0.034$: the oscillation
lasts as long as the entry.

%--------------------------------------------------------------------------
\subsection{The residual constant}\label{sec:residual-constant}
%--------------------------------------------------------------------------

A trajectory started level on the glide altitude is not on the glide: the
glide descends at the rate \cref{eq:glide-descent} and the trajectory
trails it by one phugoid response time, ringing with an altitude amplitude
$\abs{\dot h_{\mathrm{g}}}/\omega$. In scale heights that amplitude is
$C_{\calR}\sqrt{\eps}$ with
\begin{equation}\label{eq:CR}
  C_{\calR} \;=\; \frac{2\sqrt{g R_e}}{V E_v\sqrt{1 - V^2/(gr)}} ,
  \qquad\text{so that}\qquad
  C_{\calR}\sqrt{\eps} \;=\; 2\,\zeta .
\end{equation}
The residual constant and the damping ratio are one number, read two ways.

The form has been measured rather than assumed. With the scale height
varied from $2$ to $30$\,km at fixed anchor density, the largest altitude
deviation in scale heights scales as $\eps^{0.49}$ and the largest
flight-path angle as $\eps^{0.98}$, so the altitude channel sets the bound.
Against direct integration \cref{eq:CR} is an upper bound in every case
tried, tight to about four percent. The cases vary one quantity at a time
about a generic lifting body at an angle of attack $\alpha = 16^\circ$,
$V = 6$\,km/s and zero bank, the angle of attack setting $L/D$:
\begin{center}
\begin{tabular}{@{}lcccc@{}}
\toprule
case & $L/D$ & computed & \cref{eq:CR} & ratio \\
\midrule
$\alpha = 10^\circ$ & $4.89$ & $0.813$ & $0.822$ & $0.988$ \\
$\alpha = 16^\circ$ & $3.30$ & $1.193$ & $1.215$ & $0.982$ \\
$\alpha = 25^\circ$ & $2.09$ & $1.858$ & $1.912$ & $0.972$ \\
$\alpha = 35^\circ$ & $1.41$ & $2.714$ & $2.834$ & $0.958$ \\
$V = 5.0$\,km/s & $3.30$ & $1.207$ & $1.236$ & $0.977$ \\
$V = 6.8$\,km/s & $3.30$ & $1.342$ & $1.362$ & $0.985$ \\
$\sigma = 30^\circ$ & $3.30$ & $1.373$ & $1.402$ & $0.979$ \\
$\sigma = 60^\circ$ & $3.30$ & $2.335$ & $2.422$ & $0.964$ \\
\bottomrule
\end{tabular}
\end{center}
The formula bounds rather than equals because the trajectory rings: it
tracks the envelope of the oscillation, and the computed supremum touches it
from below.

%--------------------------------------------------------------------------
\subsection{Why this is not a settling rate}\label{sec:residual-settling}
%--------------------------------------------------------------------------

\Cref{eq:CR} bounds the distance to the glide for a trajectory that starts
on it. It does not say that other trajectories approach it, and by
\cref{eq:zeta} they do so only slowly. The number of oscillations per
$e$-folding of the amplitude is $1/(2\pi\zeta)$, which grows like
$\eps^{-1/2}$. In the limit the motion normal to the glide is an undamped
oscillation on its own timescale. The glide balance is then not a normally
hyperbolic slow manifold, the hypothesis under which a slow manifold
persists and attracts~\cite{fenichel1979geometric,jones1995geometric}, and
the regular-perturbation analysis of the glide notes the same thing from
the other side: the exact solution does not approach the zeroth-order
one~\cite{lu2006asymptotic}.

The frozen rate is, moreover, not what either amplitude does. The
coefficients drift as the vehicle slows, and with that drift retained the
two amplitudes of the oscillation decay at different rates,
\begin{equation}\label{eq:decay}
  \frac{\dot A_h}{A_h} = -\frac32\,\varsigma ,
  \qquad
  \frac{\dot A_\gamma}{A_\gamma} = -\frac12\,\varsigma ,
\end{equation}
for the altitude deviation and the flight-path angle respectively: the
frequency rises as the speed falls, so the altitude swing closes faster and
the angle swing slower. Their product, the action of the oscillator, decays
at $2\varsigma$. Integration of the full equations confirms both rates to
about one percent at the terrestrial scale height.

What excites the oscillation is at least as important as what damps it. A
change of bank from $\sigma_1$ to $\sigma_2$ moves the glide altitude by
$H_s\ln(\cos\sigma_2/\cos\sigma_1)$, a fixed fraction of a scale height
however thin the layer. Under bank modulation the altitude therefore
departs from the instantaneous glide by an amount of order one in scale
heights, not of order $\sqrt{\eps}$, and so does a trajectory leaving a
pass.

Two ways of closing the bound are consistent with this.
\begin{enumerate}[leftmargin=*,itemsep=2pt]
  \item \emph{Carry the amplitude as an error.} The radius of the residual
    ball is the lag $C_{\calR}\sqrt{\eps}$ plus the ringing amplitude,
    bounded at its excitation and decaying by \cref{eq:decay}. The bound
    stays explicit and the limit system stays the glide alone, at the cost
    of a residual that is of order one in scale heights after every pass
    and every bank change.
  \item \emph{Carry the amplitude as a coordinate.} The amplitude is
    promoted to a slow coordinate of the limit system and the phase is
    averaged out. The slow variables then follow the averaged inclusion and
    the fast ones converge in the statistical sense
    of~\cite{artstein1999invariant}. The residual is the averaging error,
    which for a smooth field with a single fast phase is of order
    $\sqrt{\eps}$ on the slow timescale; a rate that is uniform over the
    selections, with its constant, is what would have to be shown. The
    limit system gains a dimension.
\end{enumerate}
The body takes the second, for a reason given in
\cref{sec:residual-oscillator}: the bank history acts on the amplitude
throughout, so a bound of the first kind does not hold uniformly over bank
histories. The first remains the form the bound takes for bank histories
that vary slowly against the oscillation.

In the unstretched state variables the two descriptions agree on the
order: an oscillation of fixed amplitude in scale heights has an altitude
excursion of order $\eps R_e$ and a flight-path-angle excursion of order
$\sqrt{\eps}$, so the occupation measures do concentrate on the glide as
$\eps\to 0$, at the rate $\sqrt{\eps}$ set by the angle. What they do not
do is concentrate at that rate in the metric of \cref{eq:target}, in which
altitude is measured in scale heights.

%--------------------------------------------------------------------------
\subsection{Glide, phugoid and skip as one oscillator}
\label{sec:residual-oscillator}
%--------------------------------------------------------------------------

Let $z = (h - h_{\mathrm{g}})/H_s$ be the altitude above the glide in scale
heights. At level flight $\ddot h = V\dot\gamma = L\cos\sigma/m - G$, and
the lift at the altitude $h_{\mathrm{g}} + H_s z$ is the lift on the glide
times $e^{-z}$, that is, $G\,e^{-z}$. To leading order in $\eps$, at frozen
speed and with the damping left out, the altitude therefore obeys
\begin{equation}\label{eq:oscillator}
  \ddot z \;=\; \omega^2\bigl(e^{-z} - 1\bigr),
  \qquad \omega^2 = \frac{G}{H_s} ,
\end{equation}
which is checked as the limit $H_s\to 0$, at fixed~$z$, of the exact
vertical acceleration of the planar field. The equation is conservative,
with energy
\[
  W \;=\; \tfrac12\,\dot z^2 + \omega^2\bigl(e^{-z} + z\bigr) ,
\]
and its potential has a single minimum, at $z = 0$. Three regimes of entry
flight are three ranges of its amplitude. At the minimum the vehicle is on
the glide. Near it the potential is quadratic, with curvature $\omega^2$,
and the motion is the phugoid of \cref{eq:phugoid}. Far from it the two
sides of the well differ. Below the glide the wall is exponential, and the
trajectory rebounds from it within a few scale heights; above, the
acceleration tends to $-\omega^2$, that is, $\ddot h\to -G$, a ballistic
arc under the net gravity. A large oscillation is therefore a skip. Let
$\gamma_0$ be the flight-path angle at which the trajectory crosses the
glide altitude, so that $\dot z = V\gamma_0/H_s$ there. The rebound lasts a
time of order $H_s/(V\gamma_0)$, the arc lasts $2V\gamma_0/G$, and the two
turning points are the roots of
\begin{equation}\label{eq:amplitude}
  e^{-z} + z - 1 \;=\; \frac{V^2\gamma_0^2}{2\,G\,H_s}
  \;=\; \frac{V^2}{2\,G\,R_e}\;\frac{\gamma_0^2}{\eps} .
\end{equation}
The amplitude in scale heights is a function of $\gamma_0/\sqrt{\eps}$,
with a factor fixed by the speed. Held fixed as $\eps\to 0$, that ratio
gives an oscillation a fixed number of scale heights deep whose period is
of order $\sqrt{\eps}$: the bounded-load limit of \cref{rem:regimes}. Sent
to infinity, it gives the skip. The classical theory describes skip paths
in these terms, as oscillatory paths of large amplitude~\cite{shi1971skip};
\cref{eq:oscillator} is the equation of which that is a statement.

What the equation leaves out is of two kinds. The damping of
\cref{eq:phugoid} and the drift of the coefficients behind \cref{eq:decay}
are each of relative order $\sqrt{\eps}$ over one period. The motion of the
equilibrium when the bank changes is not small at all. With the vertical
lift scaled by $c(t) = \cos\sigma(t)/\cos\sigma_{\mathrm{ref}}$ against a
reference bank, the equation reads $\ddot z = \omega^2(c\,e^{-z} - 1)$, its
equilibrium is at $z = \ln c$, and
\begin{equation}\label{eq:pumping}
  \dot{W} \;=\; \omega^2\,(c - 1)\,e^{-z}\,\dot z .
\end{equation}
A bank history that adds lift while the vehicle climbs and removes it while
it descends raises the energy of the oscillation at every instant, and the
opposite one lowers it. The second is what a guidance law does when it
feeds the altitude rate back into the bank command, which is how the large
phugoid oscillations of vehicles of high lift-to-drag ratio are removed in
practice~\cite{lu2014entry}. The first is open to any admissible control
that can follow the oscillation, and a reachable set has to contain what it
produces. The amplitude is a controlled variable, and that is why it is
carried as a coordinate.

%--------------------------------------------------------------------------
\subsection{Range of validity}\label{sec:residual-range}
%--------------------------------------------------------------------------

$C_{\calR}$ diverges as $V^2\to gr$. Near circular speed the lift required
vanishes, the glide altitude climbs, and by \cref{eq:phugoid} the phugoid
slows until it is no longer fast: there is no separation of scales to
exploit. The constant is therefore stated on a corridor bounded away from
circular speed, and what lies above that bound is the regime of arcs and
passes, to which the glide description does not apply. The same boundary
appears independently in the attitude dynamics, where the dynamic pressure
on the glide falls with the lift required and the trim loses stiffness, so
that a floor on dynamic pressure belongs among the hypotheses of the
corridor rather than among its consequences.

%--------------------------------------------------------------------------
\subsection{The bubble budget}\label{sec:residual-budget}
%--------------------------------------------------------------------------

By \cref{lem:energy} the specific energy is non-increasing along every
trajectory of every selection, so a pass that dissipates at least $\delta$
can occur at most $K_\delta\le\Delta\mathcal{E}/\delta$ times: this is
\cref{eq:kmax}. The dissipation of a pass is tied to how far it turns the
flight path. Where the aerodynamic force dominates gravity,
$\dd V/V = -\dd\gamma/E_v$, and a pass that turns the flight path through
$\Delta\gamma$ at constant $E_v$ leaves with
\begin{equation}\label{eq:turning}
  \frac{V_{\mathrm{out}}}{V_{\mathrm{in}}}
  \;=\; e^{-\Delta\gamma/E_v} ,
\end{equation}
the classical relation for a skip, in which
$\Delta\gamma = 2\gamma_{\mathrm{entry}}$~\cite{eggers1958comparative}. A
pass entered at $4^\circ$ with $E_v = 2$ therefore dissipates about
thirteen percent of the kinetic energy, and an entry that sheds its energy
in such passes makes about fourteen of them before the speed has fallen by
a factor $e$. The count needs no heating correlation; the
integrated heat load gives the same budget its physical reading, and the
corridor constraint on it is carried separately\apxref{rem:sutton-graves-circularity}.

How deep a pass goes is fixed in the same way. Along the planar field the
quantity $\cos\gamma - E_v H_s\,\rho/(2\beta)$ changes only through
gravity, at the rate $G\sin\gamma\cos\gamma/V$: the lift and the density
gradient cancel exactly. Where the aerodynamic force dominates it is
conserved, so a pass entered from thin air at a flight-path angle
$-\gamma_e$ levels out at the density $\rho_{\mathrm{b}}$ where
\begin{equation}\label{eq:pass-depth}
  \frac{E_v H_s\,\rho_{\mathrm{b}}}{2\beta} \;=\; 1 - \cos\gamma_e ,
  \qquad\text{that is, where}\qquad
  \frac{L\cos\sigma}{m}
  \;=\; \frac{V^2\,(1-\cos\gamma_e)}{H_s}
  \;\approx\; \frac{V^2\gamma_e^2}{2H_s} .
\end{equation}
Against gravity the load at the bottom of a pass is of order
$\gamma_e^2/\eps$, and holding it fixed as the layer thins forces
$\gamma_e\sim\sqrt{\eps}$. The same identity measures what this inner
description neglects. Across a pass the conserved quantity drifts by at
most $G\gamma_e/V$ times the time spent in the layer, which is a few times
$H_s/(V\gamma_e)$, so the correction to \cref{eq:pass-depth} is of relative
order $\eps/\gamma_e^2$: the reciprocal of the same parameter, and the
small parameter of the matching in \cref{obj:composition}.

The width of a pass has been measured as well, and against what surrounds
it. A computed pass entered at $80$\,km, $6.5$\,km/s and $-4^\circ$, which
is twice $\sqrt{\eps} = 0.034$\,rad, returns to its entry altitude after
$145$\,s, having shed $14$ percent of its kinetic energy. The arc that
follows lasts $185$\,s, against the $2V\gamma_0/G = 179$\,s of
\cref{sec:residual-oscillator}, and the period $2\pi/\omega$ of the small
oscillation about the glide at the speed of the pass is $286$\,s. These are
$0.18$, $0.23$ and $0.35$ of the slow time $\sqrt{R_e/g}$. At the
terrestrial scale height the three are not separated: the pass lasts half a
period of the oscillation, and is one of its troughs more than it is an
event short against it. An estimate that substitutes
$\sqrt{\eps}\,\sqrt{R_e/g} = 27$\,s for the width of a pass understates it
fivefold.

\fi

\printbibliography[title={References}]

\end{document}
